\chapter{Raw Lecture Notes}\label{chap_Notes}

%%%%%%%%%%%%%%%%%%%%%%%%%%%%%%%%%%%%%%%%%
\section{Mini-Style Guide}
%%%%%%%%%%%%%%%%%%%%%%%%%%%%%%%%%%%%%%%%%
% Started by Arthur, please add/edit.

%%%%%%%%%%%%%%%
\begin{table}[htp] \centering
  \caption{Cheat Sheet for Symbols}
\begin{tabular}{@{\extracolsep{5pt}}l l l }\label{tbl_MiniStyleGuide}
\\[-1.8ex]\hline
\hline \\[-1.8ex]
Description & For & Type \\
\hline \\[-1.8ex]
\multirow{4}{*}{Simple symbols} & `` '' & \verb+`` ''+ \\
 & \_ & \verb+\_+ \\
 & \& & \verb+\&+ \\
 & $\sim$ & \verb+$\sim$+ \\
Hat on single character & $\hat{a}$ & \verb+$\hat{a}$+ \\
Hat on multiple chars & $\widehat{abc}$ & \verb+$\widehat{abc}$+ \\
Expectation E & $\mathbb{E}$ & \verb+$\mathbb{E}$+ \\
Scalar variable & $a$ & \verb+$a$+ \\
\multirow{2}{*}{Vector/matrix variables} & \multirow{2}{*}{$\boldsymbol{b} + \boldsymbol{\beta} + \boldsymbol{P_X}$} & \verb=$\boldsymbol{b} + \boldsymbol{\beta}= \\
 & & \verb=+ \boldsymbol{P_X}$= \\
\hline \\[-1.8ex]
\end{tabular}
\end{table}
%%%%%%%%%%%%%%%%

\subsection{Equations}
\paragraph{Equation(s) on separate line(s)} For:
\begin{eqnarray}
a=\boldsymbol{b} \qquad \text{with} \qquad 
 \boldsymbol{b}\sim\mathcal{N}(0,\sigma^{2}\boldsymbol{I}).
\end{eqnarray}
\noindent with a sentence continuing after.
\begin{verbatim}
\begin{eqnarray}
a=\boldsymbol{b} \qquad \text{with} \qquad 
 \boldsymbol{b}\sim\mathcal{N}(0,\sigma^{2}\boldsymbol{I}).
\end{eqnarray}
\noindent with a sentence continuing after.
\end{verbatim}

\subsection{Paragraphs and lists}
\paragraph{Paragraph with title} For a title like this paragraph has (especially handy for headings that are children of a \verb+section+), type:
\begin{verbatim}
\paragraph{Paragraph title}
First sentence of paragraph....
\end{verbatim}

\paragraph{Bulleted lists} For:
\begin{itemize}
\item First item.
\item Second item, etc.
\end{itemize}
Type:
\begin{verbatim}
\begin{itemize}
\item First item.
\item Second item, etc.
\end{itemize}
\end{verbatim}

\paragraph{Numbered lists} For:

\begin{enumerate}[(i)]
\item First item
\item Second item, etc. 
\end{enumerate}

Type:
\begin{verbatim}
\begin{enumerate}[(i)]
\item First item
\item Second item, etc. 
\end{enumerate}
\end{verbatim}

\subsection{Graphics}

Good introductory TikZ resources:
\begin{itemize}
\item http://www.sharelatex.com/learn/TikZ\_package
\end{itemize}

For TikZ, Arthur finds it helpful to re-sketch the graphic on graph paper, with each grid cell = 1em
\begin{itemize}
\item Regular arrowhead: \verb+-latex'+
\item Arrows/lines: with separate \verb+\path+ commands (instead of \verb+edge+ inside nodes) 
\item 3 or more objects with consistent formatting: Use \verb+\tikzstyle{}+ (instead of applying settings to each object)
\item Size specifications: use \verb+em+ (height of one line of text)
\item Positioning: use \verb+left/above/right/below of+ where possible (instead of \verb+(x,y)+ coordinates)
\end{itemize}

\subsection{General reference}
\begin{itemize}
\item http://www.giss.nasa.gov/tools/latex/ltx-2.html
\end{itemize}
%%%%%%%%%%%%%%%%%%%%%%%%%%%%%%%%%%%%%%%%%

\section{Carpenter}
%Random Variables p.1 BEGIN

\subsection{Random Variables}
\begin{enumerate}[1)]
\item Discrete random variables and expectations
categorical/qualitative:
\begin{itemize}
\item normal
\item ordinal
\item dichotomous
\end{itemize}
\item Expectation of discrete random variable\\
Expectations Operator:
\begin{eqnarray}
E(x)= x_{1}p_{1} + x_{2}p{2}+...+x_{n}p_{n}=\sum_{i=1}^{n} x_{i} p_{i}
\end{eqnarray}
Where:
\begin{eqnarray*}
\text{E(x) is the population mean}
\end{eqnarray*}
Recall that:
\begin{eqnarray*}
\sum p_{i} =1
\end{eqnarray*}
\end {enumerate}
\paragraph{Example:}
\underline{St. Petersburg Paradox:}(\textit{Bernoulli}, a scholar in fluid dynmics regarding the speed-pressure relationship)\\
Proposal: To pay $\$2^{n-1}$ if flip is tails at nth toss\\
\begin{eqnarray*}
E(x)= 1/2*\$1+1/2*1/2*\$2+1/2*1/2*1/2*\$4...
\end{eqnarray*}
\begin{eqnarray*}
E(x)= 1/2*\$1+1/4*\$2+1/8*\$4...
\end{eqnarray*}
\begin{eqnarray*}
E(x)=\sum_{i=1}^{ \infty} \$_{i}*1/2=\infty
\end{eqnarray*}
\paragraph{Example:}
\begin{eqnarray*}
E(x)= 1/6*1+1/6*2+...1/6*6=3.5
\end{eqnarray*}
With:
\begin{eqnarray*}
\mu_{DIE}=3.5 
\end{eqnarray*}
and a pair of dice 
\begin{eqnarray*}
\mu_{x}=7
\end{eqnarray*}
Generally, we can write:
\begin{eqnarray*}
E(x)=E[\delta(x)]
\end{eqnarray*}
Where E(x) is the \textit{Expectations Operator}
\begin{eqnarray*}
E(x)=\delta(x_{1}p_{i}+...\delta(x_{n})p_{n}
\end{eqnarray*}

\subsection{Expectations Rules}
\begin{enumerate}[(1)]
\item E[X+Y+Z]=E[X]+E[Y]+E[Z]
\item E[bX]=bE[X]
\item E[b]=b
\end{enumerate}
\paragraph{Example:}
\begin{eqnarray*}
E[a+bX]=a+bE(X)
\end{eqnarray*}

\subsection{Variance of a Discrete Variable}
\begin{eqnarray}
Var(x)=\delta_{x}^{2}=E[(x-\mu_{x})^{2}]\\
=(x_{1}-\mu_{x})^{2}p_{i}+(x_{2}-\mu_{x})^{2}p_{2}+...\\
=\sum_{i=1}^{n}(x_{i}-\mu_{x})^{2}p_{i}
\end{eqnarray}
\paragraph{Exercise:}
\begin{eqnarray*}
\sigma_{x}^{2}=E[(x-\mu_{x}^{2}]\\
=E[(x-\mu_{x})(x-\mu_{x})]\\
=E[x^{2}-2\mu_{x}x+\mu_{x}^2]\\
=E[x^{2}-2\mu_{x}x+\mu_{x}^2]\\
=E[x^{2}]+E[-2\mu_{x}x]=E[\mu_{x}^{2}]\\
=E[x^{2}]-2\mu_{x}E[x]+\mu_{x}^{2}\\
=E[x^{2}]-2\mu_{x}^{2}+\mu_{x}^{2}\\
=E[x^{2}]-\mu_{x}^{2}
\end{eqnarray*}

\subsection{Fixed and Random Components of Variables}
\paragraph{}
\noindent$x=\mu_{x}+u$ eg. measurement error, common trends
\begin{eqnarray}
E(u)=E[x-\mu_{x}]=E[x]-\mu_x=0
\end{eqnarray}
$\therefore$ because all variation in $x$ is due to $u$, 
\begin{eqnarray}
\delta_{x}^{2}=\delta_{u}^{2}
\end{eqnarray}
\text{Since:}
\begin{eqnarray}
\delta_{x}^{2}=E[(x-\mu_{x})^{2}]
\end{eqnarray}
\begin{eqnarray}
=E[x^{2}]-\mu_{x}^{2}
\end{eqnarray}
\begin{eqnarray}
=E[u^{2}]
\end{eqnarray}
\begin{eqnarray}
\delta_{u}^{x}=E[(u-\mu_{u})^{2})
\end{eqnarray}
\begin{eqnarray}
=E[u^{2}]
\end{eqnarray}

\subsection{Continuos Random Variables:}
\begin{eqnarray*}
\text{PDF with}f(x) \equiv F'(x)\text{CDF}
\end{eqnarray*}

%Add Figs%
%add PDF%
\begin{figure}[htp]
\caption{PDF with f(x)$\equiv$ F'(x) }\label{fig_Diversification}
\begin{center}
\includegraphics[width=0.65\textwidth]{"Figures/PDFjc"}
\end{center}
\footnotesize
\end{figure}
%add CDF%
\begin{figure}[htp]
\caption{CDF with f(x)$\equiv$ F'(x) }\label{fig_Diversification}
\begin{center}
\includegraphics[width=0.65\textwidth]{"Figures/CDFjc"}
\end{center}
\footnotesize
\end{figure}

\begin{eqnarray*}
\text{Pr.}[A\leq x\leq B]=\int_{A}^{B} \mathrm{f}(x)\,\mathrm{d}x
\end{eqnarray*}
$\therefore$
\begin{eqnarray*}
\text{CDF}=\int_{-\infty}^{x} \mathrm{f}(x)\,\mathrm{d}x=F(B)-F(A)==\int_{-\infty}^{x} \mathrm{F'}(x)\,\mathrm{d}x
\end{eqnarray*}
\text{So, the expected value of continuos random variable:}
\begin{eqnarray}
E[x]=\int x \mathrm{f}(x)\,\mathrm{d}x=\mu_{x}
\end{eqnarray}
\text{similarly}
\begin{eqnarray}
E[(x-\mu_{x})^{2}=\delta_{x}^{2}=\int(x-\mu_{x})^{2}\mathrm{f}(x)\,\mathrm{d}x
\end{eqnarray}

\paragraph{Moments of a Distribution:}
Expectations of a random variable raised to a certain power is the moment of a distribution. *Raised to the 3rd power is $\emph{Skewness}$, raised to the 4th power is $\emph{Kurtosis}$.\\ 

\begin{eqnarray}
\emph{m}_{k}(X)=\int x^{k} \mathrm{f}(x)\,\mathrm{d}x\rightarrow \mu_{k}\equiv E[X-E(x)]^{k}=\int (X-\mu_{x})^{k} \mathrm{f}(x)\,\mathrm{d}x
\end{eqnarray}

\subsection{Covariance and Correlations}
\begin{eqnarray}
Cov(X,Y)=\sigma_{xy}=E[(X-\mu_{X}(Y-\mu_{Y})]
\end{eqnarray}
$\emph{Independence}$ between X $\&$ Y if:
\begin{eqnarray}
E[\delta(X)\textit{h}(Y)]=E[\delta(X)]E[\textit{h}(Y)] 
\end{eqnarray}
$\therefore$
\begin{eqnarray}
\sigma_{XY}=0\\
\end{eqnarray}
\text{because recall that}
\begin{eqnarray}
E[X-\mu_{X}]=E[X]-\mu_{X}=0
\end{eqnarray}


\paragraph{Covariance rules:}
\begin{enumerate}[(1)]
\item Y=W+V; Cov(X,Y)=Cov(X,W)+Cov(X,V)
\item Y=b2; Cov(X,Y)=bCov(X,Z)
\item Y=b; Cov(X,Y)=0
\end{enumerate}

\paragraph{Varaince rules:}
\begin{enumerate}[(1)]
\item Y=W+V; Var(Y)=Var(W)+Var(V)+2Cov(V,W)
\item Y=b2; Var(Y)=$b^{2}$Var(Z)
\item Y=b; Var(Y)=0
\item Y=V+b; Var(Y)+Var(V)
\end{enumerate}
\paragraph{Example:}
Sample mean($\bar{x}$) on estimator of population mean ($\mu_x$)
\begin{eqnarray*}
\bar{x}=\frac{1}{n}(x_{1}+x_{2}+...+X_{n})=\frac{1}{n}\sum_{i=1}^{n}x_{i}\equiv \mu_{x}
\end{eqnarray*}

%Add Figs%
%add Esimator_1%
\begin{figure}[htp]
\caption{PDFs of x and $\bar{x}$ }\label{rawnotes_dualestimators}
\begin{center}
\includegraphics[width=0.65\textwidth]{"Figures/rawnotes_dualestimators"}
\end{center}
\footnotesize
\end{figure}

\text{$\sigma^{2}_{x}=\frac{\sigma^{2}_{x}}{n}$ due to dual structures of variables with:}
\begin{enumerate}[(i)]
\item Potential distribution of observations in sample.
\item Sample mean before and after realisation of sample.
\end{enumerate}

\begin{eqnarray}
\sigma^{2}_{x}=Var(\frac{1}{n}(x_{1}+x_{2}+...x_{n}))=
\end{eqnarray}
\begin{eqnarray}
=\frac{1}{n^{2}}(Var(x_{1})+Var(x_{2}+Var(x_{n}))
\end{eqnarray}
\begin{eqnarray}
=\frac{1}{n^{2}}(\sigma^{2}_{x}+\sigma^{2}_{x}+...+\sigma^{2}_{x})=\frac{\sigma^{2}_{x}}{n}
\end{eqnarray}
\text{$\rightarrow$ OLS estimator provides estimate of conditional mean of Y given the values of X; i.e:}
\begin{eqnarray}
E[Y|X]=\beta_{0}+\beta_{1}X
\end{eqnarray}

\subsection{Estimator Properties}
\text{(finite sample vs. asymptotic properties)}
\paragraph{Unbiasedness:}
Estimate not sytematically different from true value.
\begin{eqnarray}
E(\bar{x})=E[\frac{1}{n}(x_{1}+x_{2}+...+X_{n})]
\end{eqnarray}
\begin{eqnarray}
=\frac{1}{n}E[x_{1}+x_{2}+...x_{n}]=\frac{1}{n}(\mu_{x}+\mu_{x}+...+\mu_{x})
\end{eqnarray}
\begin{eqnarray}
=\mu_{x}
\end{eqnarray}
\paragraph{Efficiency:}
estimate produced is "close as possible to time value" (i.e. minimal dispersion around, mean)

%Add Figs%
%add Esimator_2%

\begin{figure}[htp]
\caption{PDF estimators }\label{fig_Diversification}
\begin{center}
\includegraphics[width=0.65\textwidth]{"Figures/rawnotes_estimatorab"}
\end{center}
\footnotesize
\end{figure}

$\rightarrow$ can show that sample variance of $\emph{mean}$ is the smallest of its kind.\\
$\therefore$ of unbiased estimator for mean, sample average is also efficient.\\
$\rightarrow$ efficiency and bias are relative concepts that induce a $\emph{possible trade-off}$ .\\
How do we choose either A or B? pick a function, such as MSE:

\begin{eqnarray}
\text{MSE of estimator}=E[(\hat{\theta}-\theta)^{2}]^{*}
\end{eqnarray}
=Variance of estimator + Bias$^{2}$
\paragraph{Consistency} $\hat{\theta}$ is a consistent estimator of $\theta$ if $\hat{\theta} \rightarrow \theta$ as $\emph{n}\rightarrow\infty$.
\begin{eqnarray*}
^{*}E[(\hat{\theta}-\mu\hat{\theta}-\theta)^{2}] =
\end{eqnarray*}
\begin{eqnarray*}
=E[(\hat{\theta}-\mu\hat{\theta})^{2}-2(\hat{\theta})(\mu\hat{\theta}-\theta)+(\mu\hat{\theta}-\theta)^{2}]
\end{eqnarray*}
\begin{eqnarray*}
E[(\hat{\theta}-\mu\hat{\theta})^{2}]-2(\mu\hat{\theta}-\theta)E[\hat{\theta}-\mu\hat{\theta}]+E[(\mu\hat{\theta}-\theta)^{2}]
\end{eqnarray*}

\subsection{Estimators of sample variance, covariance, correllation}
\begin{eqnarray}
S_{x}^{2}=\frac{1}{n-1}\sum_{i=1}^{n}(X_{i}-\bar{X})^{2}\rightarrow\text{unbiasedness proof ITE, pg.79}
\end{eqnarray}
\begin{eqnarray}
S_{XY}=\frac{1}{n-1}\sum^{n}_{i=1}(X_{i}-\bar{X})(Y-\bar{Y})
\end{eqnarray}
\paragraph{Normal Distribution:}
\begin{eqnarray}
\phi=\frac{1}{\delta\sqrt{2\pi}} e^{\frac{-1}{2}{\frac{x-\mu}{\sigma}}^{2}}
\end{eqnarray}
\begin{eqnarray}
X\sim N(\mu,\sigma^{2})\text{and \emph{standardised}} N(0,1)
\end{eqnarray}

\subsection{Sample size and Estimator Properties:ASYMTOTICS}
\subsection{Hypothesis Testing}
%Add Figs%
%add rejection matrix%
\begin{figure}[htp]
\caption{PDF with f(x)$\equiv$ F'(x) }\label{fig_Diversification}
\begin{center}
\includegraphics[width=0.65\textwidth]{"Figures/rawnotes_offestimator"}
\end{center}
\footnotesize
\end{figure}
\begin{eqnarray*}
H_{0}: \mu=\mu_{0} \text{vs.} H_{1}:\mu\ne\mu_{0}
\end{eqnarray*}
If null hypothesis is true, $\bar{x}$ will be distributed N $\sim$ ($\mu_{0}, \frac{\sigma^{2}}{n}$)\\

\noindent Type I error: Reject when $H_{0}$ is actually True.\\
Type II error: Non-rejection when $H_{0}$ is actually False.\\
%Add Figs%
%add rejection region curve%
\begin{figure}[htp]
\caption{95"\%" C.I. rejection regions}\label{fig_Diversification}
\begin{center}
\includegraphics[width=0.65\textwidth]{"Figures/rawnotes_rejectionregion"}
\end{center}
\footnotesize
\end{figure}
$\emph{Note:}$ $\upsilon \delta$ covers full distribution.\\
z=$\frac{\bar{x}-\mu_{0}}{s.d.(x)}\lessgtr \pm 1.96$ (actually observed value $\hat{2})$\\
s.d is not actually known; $\sigma^{2}\rightarrow s^{2}_{\bar{x}}$ use estimate of s.d.:
\begin{eqnarray}
s_{\bar{x}}=[\frac{1}{n-1}\sum(x_{i}-\bar{x})^{2}]
\end{eqnarray}
\paragraph{P-Values:}
Go beyond simple reject/not reject duality since significance levels are $\emph{essentially arbitrary}$.\\
$\rightarrow$ marginal significance level: Greatest level for which a test based on $\hat{z}$ fails to reject the null.
\begin{eqnarray}
p(\hat{z})=2(1-\phi(|\hat{z}|))
\end{eqnarray}
reject both at $5\%$, but magnitude not directly visible. $\left\{
\begin{array}{l}
\hat{z}=2.02 \\
\hat{z}=5.77
\end{array} \right. $
%Add Figs%
%add Type II curves%
\begin{figure}[htp]
\caption{Type II Error}\label{fig_Diversification}
\begin{center}
\includegraphics[width=0.65\textwidth]{"Figures/rawnotes_type2"}
\end{center}
\footnotesize
\end{figure}
$\rightarrow P_{1}=4.34 \%, P_{2}=$0.00000079$\%$\\
$\therefore$ Rejection region $\equiv$ significance level/size of test\\
Probability of rejecting false $H_{0}=1-$prob(TypeII)$=$Power of test\\
$\rightarrow$ The closer $\mu_{0}$ and $\mu_1$, the smaller the power of the test.\\
\paragraph{Confidence Interval}[around a POINT ESTIMATE vs. INTERVAL ESTIMATE]

\begin{eqnarray}
\bar{x}-1.96\sigma_{x}\leq\mu\leq\bar{x}+1.96\sigma_{x}\rightarrow(95\%CI)
\end{eqnarray}
But since $\sigma_{\bar{x}}$ is not observed, we must use estimate and corresponding $\emph{critical value}$ from the t-distribution.

\begin{eqnarray}
t=\frac{\bar{x}-\mu_{0}}{s.e.(\bar{x})}
\end{eqnarray}
\paragraph{IMPORTANT:}Ratio of 2 random variables follows a t-distribution.\\
One-sided vs. two-sided testing usually depends on context of investigation.
\begin{enumerate}[i.]
\item One-sided test common for policy evaluation; i.e. treatment vs. effect.
\item $H_{0}:\mu_{x}=0$ economic theory that determines sign of parameters.
\end{enumerate}


\paragraph{Probability Limits:}
\begin{eqnarray}
lim_{n\rightarrow\infty}P(12-\infty|>\epsilon)\rightarrow)
\end{eqnarray}
$\therefore$
\begin{eqnarray}
plim2=\infty
\end{eqnarray}
\begin{eqnarray}
lim_{n\rightarrow\infty}P(\bar{x}-\mu_{x}>\epsilon)\rightarrow)
\end{eqnarray}
$\therefore$
\begin{eqnarray}
plim\bar{x}=\mu_{x}
\end{eqnarray}

%Add Figs%
%three curves%
\begin{figure}[htp]
\caption{Three estimators with varrying s.d.}
\begin{center}
\includegraphics[width=0.65\textwidth]{"Figures/rawnotes_triestimator"}
\end{center}
\footnotesize
\end{figure}

\paragraph{Conditions for Consistency:}
\begin{enumerate}[(1)]
\item Estimator of popuation charctersitic processes plim$\therefore$distribution collapses at spike for large n.
\item Spike is located at true value
\end{enumerate}
$\rightarrow$ Unbiasedness as $\emph{sufficient}$ condition; $\emph{not necessary}$ since even biased estimators can be consistent at large sample sizes.\\
Finite vs. Large sample properties(asymptotic properties)\\
$\rightarrow$ in an analytical context plim rules are much more powerful than E[ ] which often require inddependence and other strong assumptions.\\

\paragraph{Plim Rules}
\begin{enumerate}[(i)]
\item plim(X+Y+Z)=plim(X)+plim(Y)+plim(Z)
\item plim bX=bplim(X)
\item plimb=b
\item for Z=XY, plimZ=plimX*plimY
\item for Z=$\frac{X}{Y}$, plimZ=plim(X)(plimY)^{-1}
\item plim f(x)=f(plim x)
\end{enumerate}

\paragraph{Central Limit Theorem}
$\bar{X}$ will be normally distributed, even if X is not normally distributed, given a $\emph{large enough}$ sample size.\\
$\rightarrow$ depends on underlying distribution:
\begin{enumerate}[(i)]
\item for uniform, n$\geq$ 10 suffiecient
\item for lognormal, n$\geq$100
\end{enumerate}



%%%%%%%%%%%%%%%
%%%% Hamby %%%%
%%%%%%%%%%%%%%%
%%%%Inference %%%%
%%% begin page 1 %%%
\section{Causal Inference (Introduction)}
\subsection{Target quantities of interest: Inference vs. Summarizing}

% Define block styles
\tikzstyle{block} = [rectangle, draw, fill=blue!20, 
    text width=13em, text centered, rounded corners, minimum height=4em]
\tikzstyle{line} = [draw, -latex']
\tikzstyle{cloud} = [draw, ellipse,fill=red!20, node distance=3cm,minimum height=5em]
    
\begin{tikzpicture}[node distance = 2.5cm, auto]
    % Place nodes
    \node [block] (goal) {Goals of Empirical Research};
    \node [block, below of=goal] (inf) {Inference (using facts you know to learn facts you don't know) };
    \node [block, left of=inf, xshift=-11em, node distance=4em] (summ) {Summarizing data};
    \node [block, below of=inf] (counter) {Counterfactual Inference};
    \node [block, left of=counter, xshift=-4em, node distance=4cm] (des) {Descriptive Inference};
    \node [cloud, below of=counter, node distance=4cm] (what) {"What if" question};
    \node [cloud, left of=what, node distance=4cm] (predict) {Prediction};
    \node [cloud, right of=what, node distance=4cm] (casual) {Casual inference};
    % Draw lines
    \path [line] (goal) -- (inf);
    \path [line] (goal) -- (summ);
    \path [line] (inf) -- (counter);
    \path [line] (inf) -- (des);
    \path [line, dashed] (counter) -- (what);
    \path [line, dashed] (counter) -- (predict);
    \path [line, dashed] (counter) -- (casual);
\end{tikzpicture}

%%%%Inference %%%%
%%% end page 1 %%%
%%% begin page 2 %%%

\subsection{Equivalent Linear Regression Notation}

*Standard Version:
\begin{eqnarray} 
                \notag{ \text{systematic \quad stochastic}}\\            
                \notag{\downarrow \qquad \qquad \downarrow} \qquad  \\
y_{i} = m(\underline{x}) +  e \quad = \quad x_{i} \beta \quad + \quad e_{i} \qquad \\
            e_{i} \quad \sim \quad f_{N}(e_{i}|0, \sigma^{2}) \qquad      
\end{eqnarray}

\noindent \underline{Stochastic} (from $\sigma J \tilde{o} X \sigma \sigma \equiv$  aim, guess); process that is fundamentall non-deterministic, i.e. contains a random component.

\vspace{4 mm}

*Alternative Version:
\begin{eqnarray} 
y_{i} \quad \sim \quad f_{N}(y_{i}|\mu_{i}, \sigma^{2}) \quad \leftarrow \text{stochastic}\\
\mu_{i} \quad = \quad x_{i} \beta  \qquad \quad \quad  \leftarrow \text{systematic}
\end{eqnarray}

\clearpage

\noindent Understanding the alternative regression:
Are these histograms of y (see Figures D.9, D.10, D.11, D.12) a test for normality?

%Add Figures%

\begin{figure}
\centering
\begin{minipage}{.5\textwidth}
    \centering
    \includegraphics[width=50mm,scale=0.25]{"Figures/Rplot_u1"}
    \caption{$\mu_{1}$}
\end{minipage}%
\begin{minipage}{.5\textwidth}
    \centering
    \includegraphics[width=50mm,scale=0.25]{"Figures/Rplot_u2"}
    \caption{$\mu_{2}$}
\end{minipage}
\end{figure}

\begin{figure}
\centering
\begin{minipage}{.5\textwidth}
    \centering
    \includegraphics[width=50mm,scale=0.25]{"Figures/Rplot_u3"}
    \caption{$\mu_{3}$}
\end{minipage}%
\begin{minipage}{.5\textwidth}
    \centering
    \includegraphics[width=50mm,scale=0.25]{"Figures/Rplot_u4"}
    \caption{$\mu_{4}$}
\end{minipage}
\end{figure}

\vspace{4 mm}

%%%%Inference %%%%
%%% end page 2 %%%
%%% begin page 3 %%%

\noindent *Generalized alternative notation for most statistical models:
\begin{eqnarray} 
\mathrm{y}_{i} \quad \sim \quad f(y_{i}|\theta_{i}, \alpha) \quad \leftarrow \text{stochastic}\\
\theta_{i} \quad = \quad \gamma(\textrm{X}_{i}, \beta)  \quad \leftarrow  \text{systematic}
\end{eqnarray}

\begin{table}[htb] \centering 
  \caption{Notation Explanation} 
\footnotesize 
\begin{tabular}{@{\extracolsep{8pt}} c c} 
\\[-1.8ex]\hline 
\hline \\[-1.8ex] 
 Notation & Explanation \\ 
\hline \\[-1.8ex] 
$\textrm{y}_{i}$ & random outcome variable \\
$y_{i}$ & realization of $\textrm{Y}_{i}$ \\
$f(  )$ & probability desnity function \\
$\theta$ & a systematic feature of the density that varies over observation i \\
$\gamma(  )$ & functional form (linear, exponetial...) \\
$\textrm{X}_{i}$ & explanatory variables \\
$\beta$ & effect parameters \\
$\alpha$ & ancillary parameter (feature of density that remains constant over i) \\
\hline \\[-1.8ex] 
\normalsize 
\end{tabular} 
\end{table}

%%%%Inference %%%%
%%% end page 3 %%%
%%% begin page 4 %%%

\subsection{Forms of uncertainty}

\begin{eqnarray} 
\mathrm{y}_{i} \quad \sim \quad f(y_{i}|\theta_{i}, \alpha) \quad \leftarrow \text{stochastic}\\
\theta_{i} \quad = \quad \gamma(\textrm{X}_{i}, \beta)  \quad \leftarrow \text{systematic}
\end{eqnarray}

\noindent *\underline{Estimation uncertainty}: Lack of knowledge of $\beta$ and $\alpha$. Varies as n gets larger.

\vspace{4 mm}

\noindent *\underline{Fundamental uncertainty}: Representated by the stochastic component.  Exists no matter what the researcher does, no matter how large n is.

\vspace{4 mm}

\noindent $\hookrightarrow$ Systematic component examples:

\begin{eqnarray} 
E(Y_{i})\equiv \mu_{i} = X\beta = \beta_{0} + \beta_{1}X_{1i} + ... \beta_{k}X_{ki} \\
Pr(Y_{i}=1) \equiv \pi = 1/[1+e^{x_{i}\beta}] \\
V(Y_{i}) \equiv \sigma^{2} = e^{x_{i}\beta} 
\end{eqnarray}

\noindent \underline{NB:} "$\beta$" is an effect parameter in each, but meanings differ
[ENTER CHARTS ON BOTTOM OF PAGE 4]

%%%%Inference %%%%
%%% end page 4 %%%
%%% begin page 5 %%%

\subsection{The Experimental Ideal (*Heckan, et. al. (2010) IZA working papers)}

Perry pre-school randomized experiment (1962) - High scope early childhood education approach.
\begin{enumerate}[(i)]
\item 123 black pre-schoolers in Ypsilanti, MI
\item Receive intensive intervention pre-school, home visits
\item Follow up data until 1993 (from age 3 to age 40)
    \begin{itemize}
    \item basis for Head Start program; \underline{raw evidence*} that the program worked for females, less measurable effects for males
        \begin{itemize}
        \item $\downarrow$
        \end{itemize}
    \end{itemize}
\end{enumerate}
        % Define block styles
        \tikzstyle{block} = [rectangle, draw, fill=blue!20, 
        text width=7em, text centered, rounded corners, minimum height=1.5em]
        \tikzstyle{line} = [draw, -latex']
        \qquad \qquad \qquad
        \begin{tikzpicture}[node distance = 4cm, auto]
        % Place nodes
        \node [block] (eval) {Evaluation};
        \node [block, right of=eval] (agen) {Agenda};
        \node [block, right of=agen] (form) {Formulation};
        \node [block, below of=form, xshift=-3em, node distance=1cm] (adopt) {Adoption};
        \node [block, left of=adopt] (inter) {Intervention};
         % Draw lines
        \path [line] (eval) -- (agen);
        \path [line] (agen) -- (form);
        \path [line] (form) -- (adopt);
        \path [line] (adopt) -- (inter);
        \path [line] (inter) -- (eval);
        \end{tikzpicture}

\clearpage

\vspace{4 mm}

*\underline{Selection Problem} (from correlation to causal inference)

\begin{table}[htb] \centering 
  \caption{Harding, et. al. (JUE 2012) Forclosure \& Myth} 
\footnotesize 
\begin{tabular}{@{\extracolsep{8pt}} c c c c c c c} 
\\[-1.8ex]\hline 
\hline \\[-1.8ex] 
 & & Purchase & Sale & Hold. P. & Ann \% & Std. Err. \\ 
\hline \\[-1.8ex] 
REO        & 900   & \$126k & \$246k & 7.3 & 10.4\% & (12.8\%) \\
Convential & 5,335 & \$173k & \$232k & 5.9 & 6.5\%  & (9.5\%)  \\
\hline \\[-1.8ex] 
\normalsize 
\end{tabular} 
\end{table}

\begin{table}[htb] \centering 
  \caption{Harding, et. al. (JUE 2012) Forclosure \& Myth} 
\footnotesize 
\begin{tabular}{@{\extracolsep{8pt}} c c c c } 
\\[-1.8ex]\hline 
\hline \\[-1.8ex] 
Group & Sample & Mean Health & Std. Err \\ 
\hline \\[-1.8ex] 
Hospital    & 7,774  & 3.21 & 0.014 \\
No Hospital & 90,049 & 3.93 & 0.003 \\
\hline \\[-1.8ex] 
\normalsize 
\end{tabular} 
\end{table}

\begin{eqnarray}
\hookrightarrow t-stat = \frac{\mu_{1} - \mu_{2}}{\sqrt{\gamma^{2}_{1}/N_{1} + \gamma^{2}_{2}/N_{2}}} = 58.9
\end{eqnarray}

\noindent Specify \underline{treatment} as binary variable.

\begin{eqnarray}
D_{i} = {0,1} \quad \rightarrow \text{potential outcome, irrespetive of whether treatment occurs}
\end{eqnarray}

\noindent $\therefore$ \text{Two problems with observation data:  (i) relationship only holds with \underline{certain probability}; (ii) observe \underline{actual states only}}

%%%%Inference %%%%
%%% end page 5 %%%
%%% begin page 6 %%%

\vspace{4 mm}

\noindent Potential Outcome:
\begin{eqnarray}
Y_{1i} \quad \text{if} \quad D_{i} = 1 \\
Y_{0i} \quad \text{if} \quad D_{i} = 0
\end{eqnarray}
Casual effect $\equiv Y_{1i} - Y_{0i}$ \text{(Rubin causal model; Holland (1986))}

\begin{eqnarray}
Y_{i} = \left\{
    \begin{array}{c 1}
    Y_{1i} \quad \text{if} \quad D_{i} = 1\\
    Y_{0i} \quad \text{if} \quad D_{i} = 0\\
    \end{array}\right 
\end{eqnarray}  

\begin{eqnarray}
Y_{i} = Y_{0i} + (Y_{1i} - Y_{0i})D_{i}
\end{eqnarray} 

\noindent \hookrightarrow Where $(Y_{1i} - Y_{0i})$ is the \underline{Causal Effect.}  Both potential outcomes are never observed for the same person $\therefore$ most rely on some population distributions and compare averages.

\begin{eqnarray}
E[Y_{i}|D_{i}=1] - E[Y_{i}|D_{i}=0] =  \\
E[Y_{i1}|D_{i}=1] - E[Y_{0i}|D_{i}=1] \\
+E[Y_{0i}|D_{i}=1] - E[Y_{0i}|D_{i}=0]  \\
 = E[Y_{1i} - Y_{0i}|D_{i}=1] 
\end{eqnarray}

For the above equations:
\begin{itemize}
\item Line 1 - The observed difference in average health.  \underline{PARAMETER OF INTEREST:} average treatment effect on the treated.
\item Line 2 - Mean outcomes of participants had they \underline{not} participated
\item Line 3  - This is the Selection Bias; mean outcome of the non participants cannot be used as a proxy.
\item Line 4 - AVERAGE CAUSAL EFFECT OF TREATMENT
\end{itemize}

\noindent $\therefore$ selection bias is negative in this example and can completely mask treatment effect.

%%%%Inference %%%%
%%% end page 6 %%%
%%% begin page 7 %%%

\subsection{Random assignment soloes selection bias} Selection bias (UNREADABLE WORD) because $Y_{0i}$ depends on $D_{i}$ (dependence of $Y_{0i$} and $D_{i}$).

\noindent **$D_{i}$ is rendered independent of potential outcome due to \underline{independence} of $Y_{0i}$ and $D_{i}$). 
\begin{eqnarray}
E[Y_{i}|D_{i}=1] - E[Y_{i}|D_{i}=0] =  \\
 &= E[Y_{1i}|D_{i}=1] - E[Y_{0i}|D_{i}=0]  \\
 &= E[Y_{1i}|D_{i}=1] - E[Y_{0i}|D_{i}=1]  \\
\therefore E[Y_{1i}|D_{i}=1] - E[Y_{0i}|D_{i}=1] = \\
 &= E[Y_{1i} - Y_{0i}|D_{i}=1] - E[Y_{1i}-Y_{0i}] 
\end{eqnarray}

\noindent $\hookrightarrow$ \underline{Studying the impacts of policy}

Alternatives: 
\begin{itemize}
\item full general equilibrium analysis 
\item partial equilibrium structural research
\end{itemize}

\noindent Let $D^{*}=1$ for persons who participate in program with random assignment ($r=1$ or $r=0$). 
\newline Let $D^{*}=0$ for everybody else
\begin{eqnarray}
E[Y_{1i}|D_{i}=1] = E[Y_{1i}|r=1 \quad\text{and}\quad D^{*}=1] & \text{and} \\
E[Y_{0i}|D_{i}=1] = E[Y_{1i}|r=0 \quad\text{and}\quad  D^{*}=1] &
\end{eqnarray}

%%%%Inference %%%%
%%% end page 7 %%%
%%% begin page 8 %%%

\vspace{4 mm}

\noindent \underline{NB:} Random assignment does not remove selection bias, but balances the bias between the participant and non-participant samples

\vspace{4 mm}

\newline \noindent Assume same treatment effect for all:

\begin{eqnarray}
Y_{1i} - Y_{0i} = \beta & \text{and recall that}\\
  Y_{i} = Y_{0i} + (Y_{1i}-Y_{0i})D_{i} \\
& \nonumber \\
Y_{i} = E[Y_{0i}] + [Y_{1i}-Y_{0i}]D_{i} + [Y_{0i}-E[Y_{0i}]] \\
=\alpha+\beta D_{i} + \eta_{i} \quad \qquad \qquad \qquad  \qquad \qquad\\
& \nonumber \\
\therefore E[Y_{i}|D_{i}=1] = \alpha + \beta + E[\eta_{i}|D_{i}=1] \\
 E[Y_{i}|D_{i}=0] = \alpha + E[\eta_{i}|D_{i}=0]\\
 & \nonumber \\
 \text{$\hookrightarrow$ such that} \quad E[Y_{i}|D_{i}=1] - E[Y_{i}|D_{i}=0] = \\
 \qquad \qquad \qquad \beta + E[\eta_{i}|D_{i}=1] - E[\eta_{i}|D_{i}=0] 
\end{eqnarray}

 $\hookrightarrow$ Where $\beta$ is the treatment effect, and $E[\eta_{i}|D_{i}=1] - E[\eta_{i}|D_{i}=0]$ is the selection bias.
\newline \indent $\hookrightarrow$ The selecion bias amounts to correlation between the regression error $\eta_{i}$ and the regressor $D_{i}$ 

%%%%Inference %%%%
%%% end page 8 %%%
%%% begin page 9 %%%

\vspace{4 mm}
Since,
\begin{eqnarray}
E[\eta_{i}|D_{i}=1] - E[\eta_{i}|D_{i}=0] = \\
\qquad \qquad \qquad E[Y_{0i}|D_{i}=1] - E[Y_{0i}|D_{i}=0] 
\end{eqnarray}

\noindent $\therefore$ this correlation between $\eta_{i}$ and $D_{i}$ reflects the \underline{difference in the no-treatment potential outcomes} between those who get treated and who who don't.

\vspace{4 mm}

Notice that under randomization:
\begin{eqnarray}
E[Y_{i}|r=1 \quad \text{and} \quad D_{i}^{*}=1] = \alpha +\beta + E[\eta_{i}|D_{i}^{*}=1] \\
E[Y_{i}|r=0 \quad \text{and} \quad D_{i}^{*}=1] = \alpha + E[\eta_{i}|D_{i}^{*}=1] \qquad \\
& \nonumber \\
\therefore E[Y_{i}|r=1 \quad \text{and} \quad D_{i}^{*}=1] - E[Y_{i}|r=0 \quad \text{and} \quad D_{i}^{*}=1] = \beta
\end{eqnarray}

\noindent $\hookrightarrow$ but randomization does not guarantee that the selection problem disappears (ie. $E[\eta_{i}|D_{i}^{*}=1] \neq 0$); but with randomized data selection, bias is spread across samples such that it cancels out when calculating the treatment effect.

\vspace{4 mm}

\noindent Use \underline{control variables} to get better "fix" on \beta:
\begin{eqnarray}
Y_{i} \alpha + \beta D_{i} + X_{i}\gamma + \eta_{i} 
\end{eqnarray}
$\hookrightarrow$ where $D_{i}$ and $X_{i}$ must be as uncorrelated as possible.

%%%%Inference %%%%
%%% end page 9 %%%
%%% begin page 10 %%%

\subsection{Conditional Independence Assumption (CIA)}
\begin{enumerate}[(i)]
\item Experiments ensure that causal variable of interest is indepenedent of potential outcomes.
\item CIA provides implicit justification for \underline{causal interpretation of regression estimates}
\end{enumerate}

\noindent $\hookrightarrow$ since variety of control variables (all of which are assumed to be observable) are held constant

\vspace{4 mm}

\noindent $\therefore$ if we can \underline{"select on observables}; eg. imagine $X_{i}$ as a vector that explains neighborhood characteristics, ability, family backgrounds, etc.

\vspace{4 mm}

Without CIA, $E[Y_{i}|D_{i}=1] - E[Y_{i}|D=0]$ will contain a selection bias.  CIA means that conditional on $X_{i}$, bias diasppears:
\begin{eqnarray}
{Y_{0i}, Y_{1i}} \independent \text{(independence, given)} \quad D_{i}|X_{i}
\nonumber \\
\therefore E[Y_{i}|X_{i},D_{i} =1] - E[X_{i}|X_{i}, D_{i}=0] \\
= E[Y_{i} - X_{0i}|X_{i}]
\end{eqnarray}

%%%%Inference %%%%
%%% end page 10 %%%
%%% begin page 11 %%%

\subsection{Regression Analysis in R}

\begin{eqnarray}
\text{TSS = RSS + RSS} \\
\text{PROOF:} \nonumber \\
\sum_{i=1}^{n}(y_{i}-\bar{y})^2 =\sum^n(\^y_{i}-\bar{y}_{i})^2 + \sum^n(y_{i}-\^y_{i})^2  \\
\sum_{i=1}^{n}(y_{i}-\bar{y}+\^y_{i}-\^y_{i})^2 = \sum^n(\^y_{i}-\bar{y} + \^e_{i})^2 \\
= \sum((\^y_{i}-\bar{y})^2 + 2\^e_{i}(\^y_{i}-\bar{y})+\~{e}_{i}^2)\\
= \sum(\^y_{i}-\bar{y})^2 + \sum e_{i}^2 + 2\sum \^{e}_{i}(x'\hat{\beta} - \bar{y}) \\
= \text{ESS + RSS} \\
& \nonumber \\
R^2 = \frac{ESS}{RSS} = 1 - \frac{RSS}{TSS} & \text{because} & \frac{ESS}{TSS} = \frac{TSS-RSS}{TSS} 
\end{eqnarray}

%%%%%% End page 11 %%%%%  
%%%%%% End Inference Section %%%%%%

%%%%%%%%%%%%%%%%%%%%%%%%%%%%%%%%%%%%%%%%%
\section{Prokosch}
%%%%%%%%%%%%%%%%%%%%%%%%%%%%%%%%%%%%%%%%%

\subsection{Re: Notes p.1}
Page 1 of notes skipped because an expanded version of the same headings/logic already appears in \ref{chap_SpatialStats}. Starting with p.2 below.

%SpatialStats p.2 BEGIN
\subsection{Spatial DGP}
Consider two neighbours $i$ and $j$

\begin{eqnarray}
y_{i} = \alpha_{i} + \boldsymbol{X}_{i}\boldsymbol{\beta} + e_{i} \qquad \text{with} \qquad e_i \sim i.i.d.(0, \sigma^{2}) \\
y_{j} = \alpha_{j} + \boldsymbol{X}_{j}\boldsymbol{\beta} + e_{j} \qquad \text{with} \qquad e_j \sim i.i.d.(0, \sigma^{2})
\end{eqnarray}

\paragraph{Problem} 
With more spatial dependencies, little practical usefulness since it results in a system with \emph{more parameters than observations.}

$\therefore$ potentially $n^2 - n$ relations to be estimated

Assumptions:

\begin{itemize}
\item Observations at location $i$ depend on those at location $j$ and vice versa.
\item The DGP is simultaneous across space.
\end{itemize}

The overparameterisation problem can be solved by \emph{imposing structure} in a \emph{spatial autoregressive process}

\begin{eqnarray}\label{eqn_spatialDGP}
y_{i} = \rho \sum_{i=1}^{n} W_{ij} y_{j} + e_{i}
\end{eqnarray}

\noindent where $\rho$ is the spatial correlation coefficient and $\sum_{i=1}^{n} W_{ij} y_{j}$ is the spatial lag of $y_i$, the weighted average of neighboring observations.
%SpatialStats p.2 END

%SpatialStats p.3 BEGIN
\emph{Example:}

%%%%%%%%%%%%%%%
% Begin Graphic: Adjacent regions with spatial DGP
\tikzstyle{square} = [rectangle, draw, text centered, font=\large,
    minimum width=3em, minimum height=3em]

\begin{tikzpicture}[node distance = 3.1em, auto]
    % each node represents a region. 
    % We might consider reducing node distance to 0.
    \node[square] (R1) {R1};
    \node[square, right of=R1] (R2) {R2};
    \node[square, right of=R2] (R3) {R3};
    \node[square, right of=R3] (R4) {R4};
   
    % First-degree neighbor of R1
    \path[draw, *-latex'] (R1.north) to [bend left] node [above] {N1} (R2.north);
    
    % Second-degree neighbors of R1
    % can't draw each edge with *-latex' or else we'll end up with two origin dots.
    \path[draw, *-] (R2.south) to (R2.south);
    \path[draw, -latex'] (R2.south) to [bend left] node [below] {N2} (R1.south);
    \path[draw, -latex'] (R2.south) to [bend right] node [below] {N2} (R3.south);
\end{tikzpicture}
% End Graphic 
%%%%%%%%%%%%%%%

\noindent first-order neighbours and higher-order neighbours. N.B.: R1 is its own second-order neighbour (i.e. the neighbour of the first-order neighbour.)

In matrix notation,
\begin{eqnarray}
\boldsymbol{y} = \rho\boldsymbol{Wy} + \boldsymbol{e} \qquad \text{with} \qquad \boldsymbol{e} \sim\mathcal{N}(0,\sigma^{2}\boldsymbol{I}_n)
\end{eqnarray}

In order to construct $\boldsymbol{W}$, we start with a contiguity matrix $\boldsymbol{C}$,
\begin{eqnarray}
\boldsymbol{C} = 
\begin{bmatrix}
  0 & 1 & 0 & 0\\
  1 & 0 & 1 & 0\\
  0 & 1 & 0 & 1\\
  0 & 0 & 1 & 0
\end{bmatrix}
\end{eqnarray}
\noindent where each row corresponds to an observation $i$, and each column corresponds to a neighbour $j$. Then we create $\boldsymbol{W}$ by row-standardization, i.e. row sums normalized to one, such that
\begin{eqnarray}
\boldsymbol{W} = 
\begin{bmatrix}
  0 & 1 & 0 & 1\\
  \frac{1}{2} & 0 & \frac{1}{2} & 0\\
  0 & \frac{1}{2} & 0 & \frac{1}{2}\\
  0 & 0 & 1 & 0
\end{bmatrix}
\end{eqnarray}

Subsequently,
\begin{eqnarray}
\boldsymbol{W}\boldsymbol{y} =
\begin{bmatrix}
  0 & 1 & 0 & 1\\
  \frac{1}{2} & 0 & \frac{1}{2} & 0\\
  0 & \frac{1}{2} & 0 & \frac{1}{2}\\
  0 & 0 & 1 & 0
\end{bmatrix}
\begin{bmatrix}
  y_{1}\\
  y_{2}\\
  y_{3}\\
  y_{4}
\end{bmatrix}
=
\begin{bmatrix}
  y_{2}\\
  \frac{(y_{1} + y_{2})}{2}\\
  \frac{(y_{2} + y_{4})}{2}\\
  y_{3}
\end{bmatrix}
\end{eqnarray}

Now, $\rho$ is a scalar parameter of spatial dependence with
\begin{eqnarray}
\rho
 \left\{ 
  \begin{array}{l l}
    0 & \text{no spatial autocorrelation (SA)}\\
    > 0 & \text{positive SA (clustering)}\\
    < 0 & \text{negative SA (segregation)}
  \end{array} 
 \right. % means "don't draw anything for the large close brace"
\end{eqnarray}

\subsection{Spatial Autocorrelation DGP}
\begin{eqnarray}
\boldsymbol{y} = \rho\boldsymbol{W}\boldsymbol{y} + \boldsymbol{\iota}\alpha + \boldsymbol{e}\\
(\boldsymbol{I}_n - \rho\boldsymbol{W})\boldsymbol{y} = \boldsymbol{\iota}\alpha + \boldsymbol{e}\\
\boldsymbol{y} = (\boldsymbol{I}_n - \rho\boldsymbol{W})^{-1}\boldsymbol{\iota}\alpha + (\boldsymbol{I}_n - \rho\boldsymbol{W})^{-1}\boldsymbol{e}
\end{eqnarray}

Assuming $|\rho| < 1$, we can express the inverse as\footnote{This follows because the sum of a geometric series $\sum_{k=0}^{\infty}ar^k$ with a common ratio $|r| < 1$ can be expresed as $\frac{a}{1 - r}$.}
\begin{eqnarray}
(\boldsymbol{I}_n - \rho\boldsymbol{W})^{-1} = \boldsymbol{I}_n + \rho\boldsymbol{W} + \rho^{2}\boldsymbol{W}^{2} + \rho^{3}\boldsymbol{W}^{3} + ... \qquad \text{which implies}\\
\boldsymbol{y} = \boldsymbol{\iota}\alpha + \rho\boldsymbol{W}\boldsymbol{\iota}\alpha + \rho^2\boldsymbol{W}^2\boldsymbol{\iota}\alpha + ... + \boldsymbol{e} + \rho\boldsymbol{W}\boldsymbol{e} + \rho^2\boldsymbol{W}^2\boldsymbol{e} + ...
\end{eqnarray}
\noindent Because $\alpha$ is a scalar and $\boldsymbol{W}_{1} = \boldsymbol{\iota}$, the first part converges to
\begin{eqnarray}
(1 - \rho)^{-1}\boldsymbol{\iota}\alpha \qquad \text{(sum of geometric series)}\\
\therefore \qquad \boldsymbol{y} = \frac{1}{1-\rho}\boldsymbol{\iota}\alpha + \boldsymbol{e} + \underbrace{\rho\boldsymbol{W}\boldsymbol{e} + \rho^2\boldsymbol{W}^2\boldsymbol{e} + ...}_{\text{violation of Gauss-Markov}}
\end{eqnarray}
\noindent this is a strong violation of Gauss-Markov assumption of independence of error terms, therefore OLS will be \emph{biased and inconsistent}. Consider following spatial autoregression model:
\begin{eqnarray}
\boldsymbol{y} = \rho\boldsymbol{Wy} + \boldsymbol{X\beta} + \boldsymbol{e}\\
\boldsymbol{y} = (\boldsymbol{I}_n - \rho\boldsymbol{W})^{-1}\boldsymbol{X\beta} + (\boldsymbol{I} - \rho\boldsymbol{W})^{-1}\boldsymbol{e} \qquad \text{with} \qquad \boldsymbol{e}\sim\mathcal{N}(0,\sigma^{2}\boldsymbol{I})
\end{eqnarray}

With this model, we estimate 3 parameters, $\boldsymbol{\beta}$, $\sigma^{2}$ and $\rho$. N.B. If $\rho=0$, spatial autoregression reduces to regular OLS model:
\begin{eqnarray}
\hat{\boldsymbol{\beta}}_{SAR} = (\boldsymbol{X}'\boldsymbol{X})^{-1}\boldsymbol{X}'(\boldsymbol{I} - \hat{\rho}\boldsymbol{W})\boldsymbol{y}\\
 = (\boldsymbol{X}'\boldsymbol{X})^{-1}\boldsymbol{X}'\boldsymbol{y} - \hat{\rho}(\boldsymbol{X}'\boldsymbol{X})^{-1}\boldsymbol{X}'\boldsymbol{Wy}
 = \hat{\boldsymbol{\beta}}_{OLS} - \hat{\rho}\boldsymbol{X}'\boldsymbol{X})^{-1}\boldsymbol{X}'\boldsymbol{Wy}
\end{eqnarray}
%SpatialStats p.5 END
%SpatialStats p.6 BEGIN
\subsubsection{Global Autocorrelation: Moran's $I$}

Moran's $I$ coefficient calculates the ratio between the product of a variable and its \emph{spatial lag}, with the product  of the variable adjusted for spatial weights:
\begin{eqnarray}
I = \frac
         {\frac
               {\sum_{i}\sum_{j}w_{ij}(y_{i}-\bar{y})(y_{j}-\bar{y})}
               {\sum_{i}\sum_{j}w_{ij}}
         }
         {\frac
               {\sum_{i}{(y_{i}-\bar{y})^2}}
               {n}
         }
  = [-1, +1]
\end{eqnarray}
\noindent where the minimum value of $-1$ corresponds to perfect dispersion and the maximum value of $1$ corresponds to perfect correlation.

\paragraph{Local Moran's $I$} is Moran's $I$ as the regression coefficient of a variable on its spatial lag.
\begin{eqnarray}
I = \frac
         {(y_{i}-\bar{y})\sum_{i}\sum_{j}w_{ij}(y_{i}-\bar{y})}
         {\frac
               {\sum_{i}{(y_{i}-\bar{y})^2}}
               {n}
         }
\end{eqnarray}
\noindent ...show example of ``heat map'' of local Moran's $I$ z-scores where values are significantly different from those in neighboring areas.

\subsubsection{A brief excursion on projections}

Latitude: North-south position $\equiv$ $y$-axis \\
Longitude: East-west position $\equiv$ $x$-axis

Projection is necessary to represent a 3D globe with a 2D map. Projections always distort at least one property of the original globe.

There are three main types of projections:
\begin{enumerate}[(i)]
\item Cylindrical projection: [INSERT DIAGRAM(S) HERE]
\item Conic projections: [INSERT DIAGRAM HERE]
\item Planar projections: [INSERT DIAGRAM HERE]
\end{enumerate}

%%%%%%%%%%%%%%%
\begin{table}[htp] \centering
  \caption{\textsf{R} names of projections and what they preserve}
\begin{tabular}{@{\extracolsep{5pt}}l l l l }\label{tbl_RProjections}
\\[-1.8ex]\hline
\hline \\[-1.8ex]
Preserves & Cylindrical & Conic & Planar \\
\hline \\[-1.8ex]
Area & \texttt{cylequalarea} & \texttt{albers} & \texttt{azequalarea} \\
Angle & \texttt{mercator\footnotemark[1]} & \texttt{lambert} & \texttt{stereographic} \\
 & (\texttt{utm}, \texttt{omerc}) & & \\
Distance & -- & -- & \texttt{azequidistant} \\
\hline \\[-1.8ex]
\end{tabular}

\footnotemark[1]{Transverse and oblique projections}
% In the line above I used \footnotemark by mistake, rather than
% \footnotetext. As of 2013-08-03, it seems to work (displays OK and 
% doesn't disrupt footnote numbering in text), so I'm leaving it.
% -Arthur.
\end{table}
%%%%%%%%%%%%%%%%

[INSERT DIAGRAM: 3-dimensional polar coordinates, $\phi$ = LAT, $\theta$ = LONG]

\paragraph{Different coordinate systems}
\begin{enumerate}[(i)]
\item Cartesian coordinates
\item Polar coordinates
\end{enumerate}

%%%%%%%%%%%%%%%
% Begin Graphic: polar coordinates on a cartesian grid
% loosely following http://www.texample.net/tikz/examples/intersection-of/
\tikzstyle{axis} = [very thick, -latex']
\tikzstyle{axisshadow} = [thin, dashed]
\tikzstyle{shadowlabel} = [opacity=0, text opacity=1]

\begin{tikzpicture}

    \draw[axis] (-0.1,0) -- (5,0) node(xlabel)[right]{x};
    \draw[axis] (0,-0.1) -- (0,4) node(ylabel)[above]{y};
    
    \draw[-] (0,0) -- 
        node(rlabel)[above]{$r$} 
        (4,3);
    \fill[black] (4,3) circle (1.5pt) 
        node(Plabel)[above]{$P(x,y)$};     
    
    %TODO: add curved arrow from r to x-axis with label $\theta$ 
    
    \draw[axisshadow] (0,3) -- (4,3);
    \draw[shadowlabel] (0,0) -- 
        node(Pylabel)[left]{$r sin(\theta)$}
        (0,3);
        
    \draw[axisshadow] (4,0) -- (4,3);
    \draw[shadowlabel] (0,0) -- 
        node(Pxlabel)[below]{$r cos(\theta)$}
        (4,0);

\end{tikzpicture}
% End Graphic 
%%%%%%%%%%%%%%%
\begin{eqnarray}
\therefore r = \sqrt{y^{2}+x^{2}} \qquad \text{with} \qquad
    \theta = tan^{-1}(\frac{y}{x}) = arctan(\frac{y}{x}) \qquad \text{(depending on quadrant)}
\end{eqnarray}

\paragraph{Tissot indicatrix} $\equiv$ Characterizes distortions due to map projections.

[INSERT DIAGRAM: mapping lat/long points via projection]

%SpatialStats p.9 END
%SpatialStats p.10 BLANK

%%%%%%%%%%%%%%%%%%%%%%%%%%%%%%%%%%%%%%%%%
\section{Wang}
\subsection{Basic Linear Model}
Recall
\begin{eqnarray*}
y=m(\boldsymbol{X})+e
\end{eqnarray*}
where
\begin{eqnarray*}
m(\boldsymbol{X}) \equiv \mathbb{E}[y|\boldsymbol{X}]=\sum^{\infty}_{-\infty}yf(y|\boldsymbol{X})dx
\end{eqnarray*}
\paragraph {Key properties of the regression error}
\begin{enumerate}[(1)]
\item $\mathbb{E}[e|\boldsymbol{X}]$ = 0
\item $\mathbb{E}[e]$ = 0
\item $\mathbb{E}[e\boldsymbol{X}]$ = 0
\end{enumerate}

\paragraph{Homoskedasticity vs. heteroskedasticity}
\begin{enumerate}[(a)]
\item $\mathbb{E}[e^{2}|\boldsymbol{X}]$ = $\sigma^{2}$
\item $\mathbb{E}[e^{2}|\boldsymbol{X}]$ = $\sigma^{2}(\boldsymbol{X})$
\end{enumerate}
\noindent In equation b, conditional variance $\mathbb{E}(e^{2}|\boldsymbol{X})$ is correlated with the skedastic function $m^{2}(\boldsymbol{X})$\\
\\
Old view: Variance of e varies among observations\\
New view: Conditional variance of $\sigma^{2}(\boldsymbol{X})$ depends on $\boldsymbol{X}$
\begin{eqnarray*}
\sigma^{2}=\mathbb{E}[e^{2}]=\mathbb{E}[\mathbb{E}[e^{2}|\boldsymbol{X}]]=\mathbb{E}[\sigma^{2}(\boldsymbol{X})]
\end{eqnarray*}
\begin{itemize}
\item Standard approach focuses on conditional mean only and assumes that \underline{conditional variance is fixed} (ie. homoskedasticity).
\item Economist with one foot in bucket with boiling water, other foot in bucket of ice "feels fine on average" ignore variance at your peril.
\end{itemize}

\paragraph{Linear model}
\begin{eqnarray*}
y&=&\boldsymbol{X'\beta}+e\\
\mathbb{E}[e|\boldsymbol{X}]&=& 0\\
\mathbb{E}[e^{2}|\boldsymbol{X}]&=& \sigma^{2} \text{ or } \sigma^{2}(\boldsymbol{X})
\end{eqnarray*}
Key application: 
\begin{itemize}
\item ARCH and GARCH models in finance to model time-varying volatility.
\item Robert Eagle (1982) gets 2003 Nobel Prize.
\end{itemize}

\paragraph{Best linear predictor}
\begin{eqnarray*}
F(\beta)&=&\mathbb{E}[(\overset{\overset{1\times 1}{\downarrow}}{y_{i}}-\overset{\overset{1\times k}{\downarrow}}{\boldsymbol{X_{i}'}}\overset{\overset{k\times 1}{\downarrow}}{\boldsymbol{\beta}})^2]
\end{eqnarray*}
Dropping subscripts in following lines.
\begin{eqnarray*}
F(\beta)&=&\mathbb{E}[y^2] - 2\boldsymbol{\beta}'\mathbb{E}[\boldsymbol{Xy}] + \boldsymbol{\beta}'\mathbb{E}[\boldsymbol{XX'}]\boldsymbol{\beta}
\end{eqnarray*}
$\to$ First order condition:
\begin{eqnarray*}
\frac{\partial F(\beta)}{\partial \beta} = -2\mathbb{E}[\boldsymbol{Xy}] + 2\mathbb{E}[X\boldsymbol{X'}]\beta = 0
\end{eqnarray*}
$\to$ iff $\mathbb{E}[\boldsymbol{XX'}]$ is invertable (non-singular matrix with a positive determinent)\\
$\to$ Singular matrix $\equiv$ square matrix without an inverse, zero determinant.
\begin{eqnarray*}
\therefore 2\mathbb{E}[\boldsymbol{Xy}] &=& 2\mathbb{E}[X\boldsymbol{X'}]\beta\\
\therefore \boldsymbol{\beta} &=& \mathbb{E}[\boldsymbol{XX'}]^{-1}\mathbb{E}[\boldsymbol{Xy}]
\end{eqnarray*}
Empirical analog of expected squared error is the sample average squared error.
\begin{eqnarray*}
F_n(\boldsymbol{\beta}) &=& \frac{1}{n}\sum(\boldsymbol{y}-\boldsymbol{X'\beta})^2\\
&=& \frac{1}{n} SSE_n(\boldsymbol{\beta})\\
SSE_n(\boldsymbol{\hat{\beta}}) &=& \sum{y_i}^2 - 2\boldsymbol{\beta'}\sum\boldsymbol{X_{i}}y_{i} + \boldsymbol{\beta'}\sum\boldsymbol{{X_i}{X'_i}\beta}\\
\frac{\partial SSE_n(\boldsymbol{\hat{\beta}})}{\partial \boldsymbol{\beta}} &=& -2\sum\boldsymbol{x_{i}}y_{i} + \sum\boldsymbol{X_{i}X'_{i}\beta} = 0\\
\boldsymbol{\hat{\beta}} &=& (\sum\boldsymbol{X_{i}X'_{i}})^{-1}(\sum\boldsymbol{X_{i}}y_{i})
\end{eqnarray*}

\subsection{The model in matrix notation}
\begin{eqnarray*}
y_1 &=& \boldsymbol{X'_1\beta} + e_1\\
y_2 &=& \boldsymbol{X'_1\beta} + e_1\\
\vdots\\
y_n &=& \boldsymbol{X'_n\beta} + e_n
\end{eqnarray*}
i.e.
\begin{eqnarray*}
y_i &=& \beta_0 + \beta_1 x_{i1} + \beta_2 x_{i2} + \cdots + \beta_k x_{ik}\\
\boldsymbol{X'_i}\ &=& \begin{pmatrix}x_{i1} & x_{i2} & x_{i3} & \cdots & x_{ik}\end{pmatrix}
\end{eqnarray*}
Define
\[\boldsymbol{y}=\left( \begin{array}{c}
y_1\\
y_2\\
\vdots\\
y_n \end{array} \right), \quad \boldsymbol{X}=\left( \begin{array}{c}
\boldsymbol{X'_1}\\
\boldsymbol{X'_2}\\
\vdots\\
\boldsymbol{X'_n}\end{array} \right), \quad \boldsymbol{e}=\left( \begin{array}{c}
e_1\\
e_2\\
\vdots\\
e_n\end{array} \right)\]
\\
\begin{center}
$\therefore$ y = $\boldsymbol{X\beta}$ + e
\end{center}
\begin{eqnarray*}
\text{with}\quad\sum^n_1\overset{k\times1}{\boldsymbol{X_i}}\overset{1\times k}{\boldsymbol{X_i'}} = \overset{k\times n}{\boldsymbol{X}}\overset{n\times k}{\boldsymbol{X'}}\quad\text{and}\quad\sum^n_1\boldsymbol{X_{i}}y_i = \boldsymbol{X'y}
\end{eqnarray*}
\begin{eqnarray*}
&\therefore& \quad \boldsymbol{\hat{\beta}} = (\boldsymbol{X'X})^{-1}\boldsymbol{X'y}\\
&\text{and}& \quad \boldsymbol{\hat{e}} = \boldsymbol{y-X\hat{\beta}}, \quad\sigma^2 = n^{-1}\boldsymbol{\hat{e'}\hat{e}}
\end{eqnarray*}

\subsection{Martix review}
\begin{eqnarray*}
\boldsymbol{AA^{-1}}=\boldsymbol{A^{-1}A}=\boldsymbol{I_{k}}\quad \text{for a k}\times\text{k matrix}
\end{eqnarray*}
Matrix $\boldsymbol{A}$ is said to be \underline{nonsingular} if it is of full rank; ie. rank ($\boldsymbol{A}$) = k.
\begin{eqnarray*}
\boldsymbol{A} = \begin{bmatrix} a_{1} & a_{2} & \cdots & a_{r} \end{bmatrix}\qquad k\times r
\end{eqnarray*}
has rank j iff it has j independent columns.\\
\\
$\therefore$ nonsigularity $\equiv$ full rank (no linearly dependent columns) $\to$ $\therefore$ invertibility\\
cf. nonsigularity and multicollinearity.\\
\\
\underline{Matrix determinant} is a measure of \underline{matrix volume} for as square matrix, representing a rhomboid.
\[det(A) \neq 0 \quad \text{iff }\boldsymbol{A}\text{ is non-singular}\]
Consider \begin{eqnarray*}
\boldsymbol{A}=\begin{bmatrix}
a_{11} & a_{12}\\
a_{21} & a_{22}\end{bmatrix}
\end{eqnarray*}
$det$ $\equiv$ sum of all signed elementary product
\begin{eqnarray*}
\therefore det (\boldsymbol{A}) &=& a_{11}a_{22}-a_{12}a_{21}\\
det(\boldsymbol{A^{-1}}) &=& \frac{1}{det(\boldsymbol{A})}\\
\boldsymbol{A^{-1}} &=& \frac{1}{det(\boldsymbol{A})}adj(\boldsymbol{A})\\
&=&\frac{1}{det(\boldsymbol{A})}\begin{bmatrix}
a_{11} & -a_{12}\\
-a_{21} & a_{22}\end{bmatrix}
\end{eqnarray*}
%%%%%%%%%%%%Picture on Page 6%%%%%%%%%%%%%%%
Geometry interpretation of determinants of $det(\boldsymbol{A})$\\
\begin{tikzpicture}[
    scale=1.4,
    axis/.style={very thick, ->, >=stealth'},
    pile/.style={thick, ->, >=stealth', shorten <=1pt, shorten >=1pt},
    ]
    \draw [fill=lightgray] (0,0) -- (5,2) -- (5,0) -- cycle;
    \draw [fill=lightgray] (0,0) -- (0,5) -- (3,5) -- cycle;
    \draw [fill=lightgray] (3,5) -- (5,5) -- (5,2) -- cycle;
    \draw[axis] (0,0) -- (6,0) node(xlabel)[right]{$a_{1}$};
    \draw[axis] (0,0) -- (0,6) node(ylabel)[above]{$a_{2}$};
    \draw[axisshadow] (0,5) -- (5,5);
    \draw[axisshadow] (5,0) -- (5,5);
    \fill[black] (0,0) circle (1.5pt) 
        node(O)[below]{$O$};
    \fill[black] (5,2) circle (1.5pt) 
        node(a1)[right]{$\boldsymbol{A_{1}}$};
    \fill[black] (5,0) circle (1.5pt) 
        node(a11)[below]{$a_{11}$};
    \fill[black] (0,2) circle (1.5pt) 
        node(a12)[left]{$a_{12}$};
    \fill[black] (3,5) circle (1.5pt) 
        node(a2)[above]{$\boldsymbol{A_{2}}$};
    \fill[black] (3,0) circle (1.5pt) 
        node(a21)[below]{$a_{21}$};  
    \fill[black] (0,5) circle (1.5pt) 
        node(a22)[left]{$a_{22}$};  
    \fill[black] (8,7) circle (1.5pt); 
    \draw[pile] (0,0) -- (5,2); 
    \draw[pile] (0,0) -- (3,5);
    \draw[-] (3,5) -- (5,2);
    \draw[-] (8,7) -- (5,2);
    \draw[-] (8,7) -- (3,5);
    \node at (2,6) (an1) {$\frac{1}{2}(a_{21}a_{22})$};
    \node at (8,4.3) (an2) {$\frac{1}{2}(a_{11}-a_{21})(a_{22}-a_{12})$};
    \node at (6.2,1.3) (anwhole) {$a_{11}a_{22}$};
    \node at (2,-1) (an3) {$\frac{1}{2}(a_{11}a_{22})$};
    \node at (1,4.5) (area1) {}
   edge[->,bend left=45] (an1);
    \node at (4.5,4.5) (area2) {}
   edge[->,bend right=15] (an2);
    \node at (2,0.5) (area3) {}
   edge[->,bend left=15] (an3);
    \node at (4.91,1) (areawhole) {}
   edge[->] (anwhole);
\end{tikzpicture}
\\
The area of white triangle = $a_{11}a_{12}-\frac{1}{2}a_{11}a_{12}-\frac{1}{2}a_{21}a_{22}-\frac{1}{2}a_{11}a_{12}+\frac{1}{2}a_{11}a_{12}+\frac{1}{2}a_{21}a_{22}-\frac{1}{2}a_{12}a_{21}$ = $\frac{1}{2}a_{11}a_{22}-\frac{1}{2}a_{12}a_{21}$ = $\frac{1}{2}det(\boldsymbol{A})$\\
\\
What happens if two elementary vectors $\boldsymbol{A_{1}}$ and $\boldsymbol{A_{2}}$ are linearly dependent?\\
%%%%%%%%%%%%Picture on Page 6%%%%%%%%%%%%%%%
\paragraph{ACS tables}~\\
Margins of error $\to$ permit construction of 90\% confidence bands
\begin{eqnarray*}
S.E. = \frac{(MOE_{\text{ACS}})}{1.645}
\end{eqnarray*}
Eg. Test difference in population means
\begin{eqnarray*}
\left| \frac{\mu_{1}-\mu_{2}}{\sqrt{s_{1}^{2}+s_{2}^{2}}} \right| > Z_{\text{confidence}}
\end{eqnarray*}
$\therefore$ Standard error = $\frac{MOE}{2}$  $\to$  Margin = S.E. $\times$ Z-score\\
Postnet Zip codes\\
HEXadecimal colums(RGB range 0-255)

\subsection{OLS Continued}
Excursion: $\boldsymbol{X'X}^{-1}$ as a transforming space
\begin{eqnarray*}
\boldsymbol{\hat{\beta}} = \boldsymbol{X'X}^{-1}\boldsymbol{X'y}
\end{eqnarray*}
now assume that $(\boldsymbol{X'X})^{-1}$ = $\boldsymbol{I}$, this could imply that:
\begin{enumerate}[(i)]
\item off-diagonal element 0; i.e. data is uncorrelated
\item diagonal with ones; i.e. sum of squared variable = 1 which implies $\mu_{x}=0$ and $\sigma_{x}=\frac{1}{N}$
\end{enumerate}
\begin{eqnarray*}
\because \boldsymbol{X'X}^{-1}=\boldsymbol{I} \to \boldsymbol{X'X}=\boldsymbol{I}
\therefore \boldsymbol{X'}=\boldsymbol{X^{-1}}
\end{eqnarray*}
then
\begin{eqnarray*}
\boldsymbol{\hat{\beta}}=\boldsymbol{IX'y}=\boldsymbol{X'y}\equiv\boldsymbol{X^{-1}y} \to\text{Naive Regression Coefficient}
\end{eqnarray*}
$\therefore$ The rotating matrix thus compresses, expands, and rotates the naive regression coefficients.\\
\\
$\to$ $(\boldsymbol{X'X})^{-1}$ as \underline{corrective terms} based on \underline{scale and correlations} of original regressors.\\
$\therefore$ Similar scaling interpretation of $s^{2}(\boldsymbol{X'X})^{-1}$, thus linking the variance of the residuals to the standard errors of the regression coefficients.\\
\\
\underline{Recall:}
$\boldsymbol{\hat{\beta}}$ is also a random variable with its own $\mu_{\boldsymbol{\hat{\beta}}}$ and $\sigma^{2}_{\boldsymbol{\hat{\beta}}}$.\\
Thus it has a sampling distribution.
\begin{eqnarray*}
\mathbb{E}(\boldsymbol{\hat{\beta}}|\boldsymbol{X}) &=& \boldsymbol{\beta}\quad\text{and}\quad\mathbb{E}(\boldsymbol{\hat{\beta}}) = \boldsymbol{\beta}\\
&\downarrow&\\
\mathbb{E}(\boldsymbol{\hat{\beta}}|\boldsymbol{X}) &=& \mathbb{E}[(\boldsymbol{X'X})^{-1}\boldsymbol{X'y}|\boldsymbol{X}]\\
&=& (\boldsymbol{X'X})^{-1}\boldsymbol{X'}\mathbb{E}[\boldsymbol{y}|\boldsymbol{X}]\\
&=& (\boldsymbol{X'X})^{-1}\boldsymbol{X'X\beta}\\
&=& \boldsymbol{\beta}
\end{eqnarray*}
\noindent What about the variance of $\boldsymbol{\hat{\beta}}$, Var$(\boldsymbol{\hat{\beta}})$?
\begin{eqnarray*}
Var(\boldsymbol{e}) &=& \mathbb{E}(\overset{n\times1}{\boldsymbol{e}}\overset{1\times n}{e'}) = \overset{1\times1}{\sigma^{2}}\overset{n\times n}{\boldsymbol{I}}\quad\text{and}\\
\boldsymbol{\hat{\beta}}-\boldsymbol{\beta} &=& (\boldsymbol{X'X})^{-1}\boldsymbol{X'e}
\end{eqnarray*}
\begin{eqnarray*}
V(\boldsymbol{\hat{\beta}}) &=& \mathbb{E}[(\boldsymbol{\hat{\beta}-\beta})(\boldsymbol{\hat{\beta}-\beta})']\\
&=& \mathbb{E}[(\boldsymbol{XX'})^{-1}\boldsymbol{X'ee'X}(\boldsymbol{X'X})^{-1}]\\
&=& (\boldsymbol{X'X})^{-1}\boldsymbol{X'}\mathbb{E}(\boldsymbol{ee'})\boldsymbol{X}(\boldsymbol{X'X})^{-1}\\
&=& (\boldsymbol{X'X})^{-1}\boldsymbol{X'}\sigma^{2}\boldsymbol{I}\boldsymbol{X}(\boldsymbol{X'X})^{-1}\\
&=& \sigma^{2}(\boldsymbol{X'X})^{-1}\boldsymbol{X'X}(\boldsymbol{X'X})^{-1}\\
&=& \sigma^{2}(\boldsymbol{X'X})^{-1}
\end{eqnarray*}
\paragraph{Estimate of error variance $\sigma^{2}$} ~\\
OLS does not provide a direct estimate of $\sigma^{2}$, if $\sigma$ was actually observable, then
\begin{eqnarray*}
\sigma^{2} &=& \frac{1}{n}\sum e^{2}_{i}
\end{eqnarray*}
Why? The mean of $e_i$ is zero!
But best we can do is actually
\begin{eqnarray*}
\hat{\sigma}^{2} &=& \frac{1}{n}\sum \hat{e}^{2}_{i}
\end{eqnarray*}
Therefore, each squared residual can be shown to have an expectation that is smaller than $\sigma$ such that
\begin{eqnarray*}
\mathbb{E}[\hat{\sigma}^{2}] < \sigma^{2}\qquad\therefore\hat{\sigma}^{2}\text{\underline{always biased downwards}}
\end{eqnarray*}
We can show that:
\begin{eqnarray*}
\mathbb{E}[\hat{\sigma}^{2}] = \frac{n-k}{n}\sigma^{2}
\end{eqnarray*}
$\therefore$ Bias takes "scale form" (depends on ratio of coefficients to sample size $frac{k}{n}$), thus unbiased estimator via re-scaling.
\begin{eqnarray*}
s^{2}\equiv\frac{1}{n-k}\sum^{n} \hat{e_{i}}^{2}=\frac{n}{n-k}\hat{\sigma}^{2}
\end{eqnarray*}
Note: While $s^{2}$ is an unbiased estimate of $\sigma^{2}$, s does \underline{not} provide unbiased estimate of $\sigma$!

\subsection{Basic matrix algebra revisited}
\begin{eqnarray*}
\boldsymbol{a_{i}}=\left( \begin{array}{c}
a_{1i}\\
a_{2i}\\
\vdots\\
a_{ni} \end{array} \right) \quad\to\quad\boldsymbol{A}=[\boldsymbol{a_{1} a_{2}}\cdots \boldsymbol{a_{k}}] = \begin{bmatrix}
  \boldsymbol{\alpha_{1}} \\
  \boldsymbol{\alpha_{2}}\\
  \vdots\\
  \boldsymbol{\alpha_{n}}
\end{bmatrix} = \begin{bmatrix}
  a_{11} & a_{12} & \cdots & a_{1k} \\
  a_{21} & a_{22} & \cdots & a_{2k} \\
  \vdots & \vdots & \ddots & \vdots \\
  a_{n1} & a_{n2} & \cdots & a_{nk}
\end{bmatrix}
\end{eqnarray*}
\underline{Column vector}\\
$\boldsymbol{a^{T}_{i}}$ = $\boldsymbol{\alpha_{i}}$ = $\begin{bmatrix} \alpha_{j1} & \alpha_{j2} & \cdots & \alpha_{jk} \end{bmatrix}$\\
There is a special vector $\boldsymbol{\iota}$ with $\boldsymbol{\iota^{T}a}$ = $\sum \boldsymbol{a_{i}}$ = n$\overline{a}$ \\
The inner product for $\boldsymbol{\iota}$ and $\boldsymbol{a}$ is the summation of vector elements. $\therefore$ $\boldsymbol{\iota^{T}\iota}$ = n\\
\\
\underline{Row vector}\\
There are three ways to write a linear model. Recall the OLS model in matrix notation:\\
\\
\begin{tikzpicture} \node[draw,circle,inner sep=1pt]{1};
\end{tikzpicture}
\begin{eqnarray*}
&y_{1}& = \boldsymbol{x_{1}^{'}\beta} + e_{1}\\
&y_{2}& = \boldsymbol{x_{2}^{'}\beta} + e_{2}\\
&\vdots&\\
&y_{n}& = \boldsymbol{x_{1}^{'}\beta} + e_{n}
\end{eqnarray*}
with
\begin{eqnarray*}
\boldsymbol{y}=\left( \begin{array}{c}
y_1\\
y_2\\
\vdots\\
y_n \end{array} \right), \quad \boldsymbol{X}=\left( \begin{array}{c}
\boldsymbol{X'_1}\\
\boldsymbol{X'_2}\\
\vdots\\
\boldsymbol{X'_n}\end{array} \right), \quad \boldsymbol{e}=\left( \begin{array}{c}
e_1\\
e_2\\
\vdots\\
e_n\end{array} \right)
\end{eqnarray*}
\begin{tikzpicture} \node[draw,circle,inner sep=1pt]{2};
\end{tikzpicture}
\begin{eqnarray*}
\boldsymbol{y} = \boldsymbol{X\beta} + \boldsymbol{e} 
\end{eqnarray*}
\begin{tikzpicture} \node[draw,circle,inner sep=1pt]{3};
\end{tikzpicture}
\begin{eqnarray*}
y_{i}=\beta_{0}+\beta_{1}x_{1i}+\beta_{2}x_{2i}+\cdots+\beta_{k-1}x_{k-1,i}+e_{i}
\end{eqnarray*}

\paragraph{Important notational equivalents}
\begin{eqnarray*}
\sum^{n}_{i=1} \boldsymbol{x_{i}x_{i}^{'}} &=& \boldsymbol{X'X}\\
\sum^{n}_{i=1} \boldsymbol{x_{i}}y_{i} &=& \boldsymbol{X'Y}\\
\text{Since}\quad\boldsymbol{X}=\begin{bmatrix}
\boldsymbol{X'_1}\\
\boldsymbol{X'_2}\\
\vdots\\
\boldsymbol{X'_n}
\end{bmatrix}\quad&\text{and}&\quad\boldsymbol{X'}=\begin{bmatrix}\boldsymbol{x_{1}}&\boldsymbol{x_{2}}&\cdots&\boldsymbol{x_{n}}\end{bmatrix}
\end{eqnarray*}
\begin{eqnarray*}
\therefore\quad\boldsymbol{X'X}&=&\begin{bmatrix}\boldsymbol{x_{1}}&\boldsymbol{x_{2}}&\cdots&\boldsymbol{x_{n}}\end{bmatrix}
\begin{bmatrix}
\boldsymbol{X'_1}\\
\boldsymbol{X'_2}\\
\vdots\\
\boldsymbol{X'_n}
\end{bmatrix}\\
&=&\boldsymbol{X_{1}X_{1}^{'}}+\boldsymbol{X_{2}X_{2}^{'}}+\cdots+\boldsymbol{X_{n}X_{n}^{'}}\\
&=&\sum^{n}_{i=1}\boldsymbol{x_{i}x_{i}^{'}}\quad\to
\begin{bmatrix}
\sum \boldsymbol{x_{1}}^{2} & \sum\boldsymbol{x_{1}x_{2}} & \cdots & \sum\boldsymbol{x_{1}x_{k}}\\
\sum \boldsymbol{x_{2}x_{1}} & \sum\boldsymbol{x_{2}x_{2}} & & \sum\boldsymbol{x_{2}x_{k}}\\
\vdots & \vdots & \ddots & \vdots\\
\sum \boldsymbol{x_{k}x_{1}} & \sum\boldsymbol{x_{k}x_{2}} & & \sum\boldsymbol{x_{k}}^{2}
\end{bmatrix}
\end{eqnarray*}
Also, sum of square $\sum x_{i}^{2}=\boldsymbol{x'x}$ (for $\boldsymbol{x}$ as a column vector)

\subsection{Centering of variables (deviation from mean)}
\paragraph{Example} $\boldsymbol{x}$ and its mean $\overline{\boldsymbol{x}}=\frac{1}{n}\sum \boldsymbol{x_{i}}$\\
Let Z$\equiv\boldsymbol{x}-\overline{\boldsymbol{x}}\boldsymbol{\iota}$, we can show that the centered variable $\boldsymbol{\iota}$ is orthogonal (inner product is zero) to the constant $\xi$.
\begin{eqnarray*}
\boldsymbol{\iota^{'}Z}=\boldsymbol{\iota'}(\boldsymbol{X}-\overline{\boldsymbol{X}}\boldsymbol{\iota})=\boldsymbol{\iota'X}-\overline{\boldsymbol{X}}\boldsymbol{l'l}=n\overline{\boldsymbol{X}}-\overline{\boldsymbol{X}}n=0
\end{eqnarray*}
Became centering leads to a variable that is othogonal to $\boldsymbol{\iota}$, it can be performed by the orthogonal projection matrix $\boldsymbol{M_{\iota}}$
\begin{eqnarray}
\boldsymbol{M_{\iota}}\equiv\boldsymbol{I_{\iota}}-\boldsymbol{P_{\iota}}=\boldsymbol{I}-\boldsymbol{\iota}(\boldsymbol{\iota'\iota})^{-1}\boldsymbol{\iota'}=\boldsymbol{I}-\frac{\boldsymbol{\iota\iota'}}{n}
\end{eqnarray}
Proof:
\begin{eqnarray*}
\boldsymbol{M_{\iota}}&=&(\boldsymbol{I}-\boldsymbol{P_{\iota}})\\
&=&\boldsymbol{X}-\boldsymbol{\iota}(\boldsymbol{\iota'\iota})^{-1}\boldsymbol{\iota'X}\\
&=&\boldsymbol{X}-\overline{\boldsymbol{X}}\boldsymbol{\iota}\\
&=&\boldsymbol{Z}
\end{eqnarray*}
\paragraph{Application to regression analysis $R^2$:} from the goodless of fit to specification\\
Recall that:
\begin{eqnarray*}
TSS &=& ESS+RSS\\
\sum^{n}(y_{i}-\overline{y})^{2} &=& \sum^{n}(\hat{y_{i}}-\overline{y})^{2}+\sum^{n}(y_{i}-\hat{y_{i}})^{2}\\
 &=& \sum^{n}(\hat{y_{i}}-\overline{y})^{2}+\sum^{n}\hat{e_{i}}^{2}
\end{eqnarray*}
Plm: $\boldsymbol{y} = \boldsymbol{P_{x}y} + \boldsymbol{M_{x}y}\quad\therefore \boldsymbol{\hat{y}}=\boldsymbol{P_{x}y}\quad\to\quad \boldsymbol{y}-\overline{\boldsymbol{y}}\quad\equiv\quad$centered variable by multiplying by $\boldsymbol{M_{\iota}}$\\
$\therefore\boldsymbol{M_{\iota}}=\boldsymbol{y}-\overline{\boldsymbol{y}}\quad$ and recall that $\quad\boldsymbol{y'y}=\sum y^{2}$
\begin{eqnarray*}
\therefore TSS &=& \boldsymbol{y'M_{\iota}y}\\
ESS &=& \boldsymbol{y'M_{\iota}P_{x}y}\\
RSS &=& \boldsymbol{y'M_{x}y}\\
\text{And thus:}\quad R^{2} &=& \frac{ESS}{TSS}=1-\frac{RSS}{TSS}
\end{eqnarray*}
NB: Adjusted $R^{2}=1-(1-R^{2})\frac{n-1}{n-k}$ with k-1 regressors.\\
$\to\overline{R}^{2}$ increases only if the absolute value of t-stats of new variable is greater than 1.\\
$\therefore\quad$Increase in $\overline{R}^{2}$ generally not indicative of improved specification.\\
$\to$Applied econometricians increasingly pay less and less attention to specification optimisation via $R^{2}$. Evan a poorly specified equation can archive a higher $R^{2}$.

\subsection{Geometry of GIS}
$\hat{\beta}$ is a random variable with its own $\mu_{\hat{\beta}}$ and $\sigma^{2}_{\hat{\beta}}$
\begin{eqnarray*}
\boldsymbol{\hat{\beta}} &=& (\boldsymbol{X'X})^{-1}\boldsymbol{X'y}\\
&=& (\boldsymbol{X'X})^{-1}\boldsymbol{X'}(\boldsymbol{X\beta}+\boldsymbol{e})\\
&=& (\boldsymbol{X'X})^{-1}\boldsymbol{X'X\beta}+(\boldsymbol{X'X})^{-1}\boldsymbol{e}\\
&=& \boldsymbol{\beta}+(\boldsymbol{X'X})^{-1}\boldsymbol{X'e}
\end{eqnarray*}
\begin{eqnarray*}
\mathbb{E}(\boldsymbol{e}|\boldsymbol{X})=0\quad\to\quad\boldsymbol{\hat{y}}=\boldsymbol{X\hat{\beta}}=\boldsymbol{X}(\boldsymbol{X'X})^{-1}\boldsymbol{X'y}=\boldsymbol{P_{x}y}
\end{eqnarray*}
\underline{Orthogonal projections}\\
$\boldsymbol{P_{x}}$ and $\boldsymbol{M_{x}}$ are complementary projection matrices, $\boldsymbol{P_{x}}+\boldsymbol{M_{x}}=\boldsymbol{I}$
\begin{eqnarray*}
\boldsymbol{P_{x}} &=& \boldsymbol{X}(\boldsymbol{X'X})^{-1}\boldsymbol{X'}\\
\boldsymbol{M_{x}} &=& \boldsymbol{I}-\boldsymbol{P_{x}}=\boldsymbol{I}-\boldsymbol{X}(\boldsymbol{X'X})^{-1}\boldsymbol{X'}
\end{eqnarray*}
$\boldsymbol{P_{x}}$ and $\boldsymbol{M_{x}}$ are orthogonal complements which restore the original vector $\boldsymbol{y}=\boldsymbol{P_{x}y}+\boldsymbol{M_{x}y}$\\
$\boldsymbol{P_{x}}$ is idempotent matrix since $\boldsymbol{P_{x}P_{x}}=\boldsymbol{P_{x}}$\\
Properties of $P_{x}$ and $M_{x}$:
\begin{eqnarray*}
\boldsymbol{P_{x}M_{x}} &=& 0 \quad\text{since}\quad \boldsymbol{P_{x}}(\boldsymbol{I}-\boldsymbol{P_{x}})=\boldsymbol{P_{x}}-\boldsymbol{P_{x}P_{x}}=0\\
\boldsymbol{M_{x}X} &=& 0\\
\boldsymbol{P_{x}y} &=& \boldsymbol{X\hat{\beta}}=\boldsymbol{\hat{y}}\\
\boldsymbol{M_{x}y} &=& (\boldsymbol{I}-\boldsymbol{P_{x}})\boldsymbol{y}=\boldsymbol{y}-\boldsymbol{P_{x}y}=\boldsymbol{y}-\boldsymbol{X\hat{\beta}}=\boldsymbol{y}-\boldsymbol{\hat{y}}=\boldsymbol{\hat{e}}
\end{eqnarray*}
%Picture
\begin{tikzpicture}[
    scale=2,
    axis/.style={very thick, ->, >=stealth'},
    pile/.style={thick, ->},
    connect/.style={thick, -},
]
    \draw[-] (0,1) -- (2,5) -- (5,4) -- (3,0) -- (0,1);
    \fill[black] (1.25,1.75) circle (1pt) 
        coordinate(o);
    \fill[black] (2,1.5) circle (1pt) 
        coordinate(x2);
    \fill[black] (1.75,2.75) circle (1pt) 
        coordinate(x1);
    \fill[black] (2.5,2.5) circle (1pt) 
        coordinate(a);
    \draw[axis] (o) -- (2.75,1.25)
        node(x2a)[right]{$\boldsymbol{X_{2}}$};
    \draw[axis] (o) -- (2.70,4.65)
        node(x1a)[right]{$\boldsymbol{X_{1}}$};
    \draw[axis] (o) -- (2.5,5.89)
        node(y)[right]{$\boldsymbol{y}$};
    \draw[axis] (o) -- (a);
    \draw[axis] (o) -- (x1);
    \draw[axis] (o) -- (x2);
    \draw[axisshadow] (a) -- (x1);
    \draw[axisshadow] (a) -- (x2);
    \draw[axisshadow] (a) -- (2.5,5.89);
    \draw[-] (2.5,2.7) -- (2.3,2.58) -- (2.4,2.44);
    \node at (-0.3,2.5) (pxy) {$\boldsymbol{P_{x}y}$}
   edge[->,bend left=15] (1.5,1.9);
    \node at (4.5,3.8) (mxy) {$\boldsymbol{M_{x}y}$}
   edge[->,bend left=15] (2.5,3.5);
    \node at (0.75,1.5) (x1b1) {$\boldsymbol{X_{1}\hat{\beta_{1}}}$}
   edge[->,bend left=15] (1.5,2.25);
   \node at (2,1.25) (x2b2) {$\boldsymbol{X_{2}\hat{\beta_{2}}}$};
\end{tikzpicture}
\\
%Picture
\begin{tikzpicture}[
    scale=2,
    axis/.style={very thick, ->, >=stealth'},
    pile/.style={thick, ->},
    connect/.style={thick, -},
    ]
    \fill[black] (2,0) circle (1pt) 
        coordinate(o) node[below]{$O$};
    \fill[black] (5,0) circle (1pt) 
        coordinate(x1);
    \fill[black] (0,1) circle (1pt) 
        coordinate(x2);
    \fill[black] (2.8,2) circle (1pt) 
        coordinate(a) node[left]{$A$};
    \fill[black] (3,1) circle (1pt) 
        coordinate(b) node[right]{$B$};
    \fill[black] (2.8,0) circle (1pt) 
        coordinate(c) node[below]{$C$};  
    \draw[axis] (o) -- (-0.4,1.2) node(lx2)[left]{$\boldsymbol{X_{2}}$};
    \draw[axis] (o) -- (5.5,0) node(lx1)[right]{$\boldsymbol{X_{1}}$};
    \draw[connect] (x1) -- (b);
    \draw[connect] (x2) -- (b);
    \draw[axisshadow] (a) -- (c);
    \draw[pile] (o) -- (b);
    \draw[pile] (b) -- (a);
    \draw[connect] (c) -- (b); 
    \draw[axis] (o) -- (3,2.5)
        node(y)[above]{$\boldsymbol{y}$};
    \draw[-] (2.8,0.15) -- (2.95,0.15) -- (2.95,0);
    \draw[-] (2.6,0) -- (2.7,0.2) -- (2.8,0.15);
    \draw[-] (2.9,0.9) -- (2.8,1) -- (2.97,1.15);
    \node at (2.2,-0.8) (pxy) {$\boldsymbol{P_{x}y}$}
   edge[->,bend right=15] (2.3,0.3);
    \node at (3.5,1.8) (mxy) {$\boldsymbol{M_{x}y}$}
   edge[->,bend left=15] (2.92,1.4);
   \node at (0.52,0.5) (x2b2) {$\boldsymbol{X_{2}\hat{\beta_{2}}}$};
   \node at (4,-0.2) (x1b1) {$\boldsymbol{X_{1}\hat{\beta_{1}}}$};
\end{tikzpicture}
~\\
\noindent OLS as a more general estimator:\\
\\
$\boldsymbol{\tilde{\beta}}=(\boldsymbol{Z'Z})\boldsymbol{Z'y}\quad$ eg. $\quad\boldsymbol{\hat{\beta}}_{2SLS}\equiv\boldsymbol{\beta}_{IV}$ if just-identified.\\
\begin{center}
\begin{tikzpicture}
    \matrix (m) [matrix of math nodes,row sep=1.5em,column sep=1.5em,minimum width=1em]
  {
    x & y & \text{but what if} & \overset{\text{ability}}{z} & \overset{\text{schooling}}{x} & \overset{\text{wage}}{y} \\
    e & & & & e & \\};
  \path[-stealth]
    (m-1-1) edge (m-1-2)
    (m-2-1) edge (m-1-2)
    (m-1-4) edge (m-1-5)
    (m-1-5) edge (m-1-6)
    (m-2-5) edge (m-1-5)
    (m-2-5) edge (m-1-6);
\end{tikzpicture}
\end{center}
%This should use {align} instead of {eqnarray}%
\begin{eqnarray*}
\boldsymbol{\hat{\beta}}_{2SLS} = [\boldsymbol{X'Z}(\boldsymbol{Z'Z})^{-1}\boldsymbol{Z'X}]^{-1}[\boldsymbol{X'Z}(\boldsymbol{Z'Z})^{-1}\boldsymbol{Z'y}]
\end{eqnarray*}
With first stage\\
\begin{eqnarray*}
\boldsymbol{\hat{X}} &=& \boldsymbol{Z}\boldsymbol{Z'Z})^{-1}\boldsymbol{Z'X}\\
\text{i.e.}\quad\boldsymbol{\hat{\beta}}_{\text{1stage}} &=& (\boldsymbol{Z'Z})^{-1}\boldsymbol{Z'X}\qquad\boldsymbol{\hat{\beta}}_{V}=(\boldsymbol{Z'X})^{-1}\boldsymbol{X'y}
\end{eqnarray*}

\subsection{GDP and GDI}
\paragraph{\underline{GDI}(2003)}
\begin{itemize}
\item Compensation of employees (56\%)
\item Taxes on production and imports (TOPI) - subsides on production and imports (7\%)
\item Net operating surplus (25\%)
\item Comsumption of fix capital (CFC) (12\%)
\end{itemize}
GDI + Statistical discrepancy (about 0.2\%) $\equiv$ GDP
\begin{center}
\begin{tabular}{rrp{8cm}}
TOPI & $\equiv$ & sales taxes, property taxes, custom duties\\
NOS & $\equiv$ & corporate profits, net exchange of unincorporated business, government surplus. eg. FNMA, federal profits\\
CFC & $\equiv$ & measured depreciation, capital expenditure, interest
\end{tabular}
\end{center}
\paragraph{\underline{Measuring price changes}}
\begin{itemize}
\item CPI: fixed bucket of goods over time
\item GDP deflator: changes in nominal GDP over time ($\frac{\text{nominal GDP}}{\text{real GDP}}\times 100$)
\end{itemize}
$\to$ Main difference to CPI is that weights of bucket is not fixed.
\paragraph{Two types of \underline{fixed-weight indices}}
\begin{enumerate}[(1)]
\item Paasche index: $\frac{\sum p_{it}q_{it}}{\sum p_{i0}q_{it}}$, understates inflation
\item Laspeyes index: $\frac{\sum p_{it}q_{i0}}{\sum p_{i0}q_{i0}}$, overstates inflation
\end{enumerate}
$\to$ CPI-W, CPI-U until recently as Laspeyes indices. Now a Fisher ideal indices which is \underline{geometrix mean} between Paasche and Laspeyes.
\begin{eqnarray*}
(\prod^{n}_{i=1}a_{i})^{\frac{1}{n}}\quad\text{vs.}\quad\frac{1}{n}\sum^{n}_{i=1}a_{i}
\end{eqnarray*}

\subsection{Transformation of variables (Nonlinearity)}
General assumptions of classical linear regression model
\begin{itemize}
\item dependent variables as linear function of independent variables
\item $\beta$ is constant
\end{itemize}
which may leads to $\to$ \begin{enumerate}[(i)]
\item wrong regression (inclussion or omission)
\item nonlinearity
\item changing parameters ($\beta$ varies with time or actual observation. eg. random coefficient model)
\end{enumerate}
\begin{tikzpicture} \node[draw,circle,inner sep=0.8pt]{1};
\end{tikzpicture} Transform single dependent variables
\begin{eqnarray*}
y &=& \beta_{0}+\beta_{1}x+\beta_{2}x^{2}+e\\
\to \text{let}\quad z &=& x^{2} \quad \text{to preserve linearity in parameters}
\end{eqnarray*}
\begin{tikzpicture} \node[draw,circle,inner sep=0.8pt]{2};
\end{tikzpicture} Transform entire equations
\begin{eqnarray*}
\text{Cobb-Donglas production function}\quad Y=AK^{\alpha}L^{\beta}e \quad\text{with}\quad CES\quad \alpha=1-\beta\\
\text{with log linearisation:}\quad \ln Y=\ln A + (1-\beta)\ln K + \beta\ln L + \ln e
\end{eqnarray*}
\paragraph{Constant return to scale}
\begin{eqnarray*}
F(aK,aL) &=& A(aK)^{\alpha}(aL)^{1-\alpha}\\
&=& Aa^{\alpha}K^{\alpha}a^{1-\alpha}L^{1-\alpha}\\
&=& a(AK^{\alpha}L^{1-\alpha})
\end{eqnarray*}
What if ($\alpha + \beta$) $<$ 1?
\paragraph{Elasticity (of substitution):}
\begin{eqnarray*}
\frac{dY}{Y}\cdot\frac{X}{dX} &=& \frac{\frac{dY}{dX}}{\frac{Y}{X}} \equiv \frac{\text{marginal product}}{\text{average product}}\\
\frac{dT}{dL} &=& \beta AK^{\alpha}L^{\beta-1} = \beta\frac{T}{L}\\
\therefore \frac{\frac{dY}{dL}}{\frac{Y}{L}} &=& \beta\frac{T}{L}\cdot\frac{L}{Y} = \beta\quad\text{[Constant factor proportion]}
\end{eqnarray*}
Cobb-Donglas functions derived from CES:
\begin{eqnarray*}
Y=A[\alpha K^{\theta}+(1-\alpha)L^{\theta}]^{\frac{1}{\theta}} \\
\theta=0\text{ yields}\quad Y=AK^{\alpha}L^{1-\alpha}\quad\text{from}
\end{eqnarray*}
\begin{eqnarray*}
\ln(Y) &=& \ln(A)+\frac{1}{\theta}\cdot\ln[\alpha K^{\theta}+(1-\alpha L^{\theta})]\\
\text{L'Hospital's rule:}\quad\lim_{\theta\to 0}\ln(Y) &=& \ln(A)+\alpha\ln(K)+(1-\alpha)\ln(L)\\
\therefore\quad Y &=& AK^{\alpha}L^{1-\alpha}
\end{eqnarray*}
\paragraph{Frisch-Waugh-Lovell Theorem (FWL)} ~\\
if $\boldsymbol{Y}=\boldsymbol{X_{1}}\beta_{1}+\boldsymbol{X_{2}}\beta_{2}+\boldsymbol{e}$ where $\boldsymbol{X_{1}}$ and $\boldsymbol{X_{2}}$ are $n \times k_{1}$, and $n \times k_{2}$ matrices, and where $\boldsymbol{\beta_{1}}$ and $\boldsymbol{\beta_{2}}$ are conformable then the estimate of $\boldsymbol{\beta_{2}}$ will be the same as the estimate from an anxilliary regression of the form
\begin{eqnarray}
\boldsymbol{M_{X_{1}}y}=\boldsymbol{M_{X_{1}}X_{2}\beta_{2}}+\boldsymbol{\upsilon}\quad\text{(where }\boldsymbol{\upsilon}=\boldsymbol{M_{X_{1}}e}\text{)}
\end{eqnarray}
where $\boldsymbol{M_{X}}$ projects onto the orthogonal complement of the image of the projection matrix $\boldsymbol{X_{1}}(\boldsymbol{X'_{1}X_{1}})^{-1}\boldsymbol{X'_{1}}$.\\
\\
Equivalently, $\boldsymbol{M_{X}}$ projects onto the orthogonal complement of the column space of $\boldsymbol{X_{1}}$. Specially,
\begin{eqnarray}
\boldsymbol{M_{X_{1}}}=\boldsymbol{I}-\boldsymbol{X_{1}}(\boldsymbol{X'_{1}X_{1}})^{-1}\boldsymbol{X'_{1}}
\end{eqnarray}
Proof:
\begin{eqnarray*}
\text{Let}\quad\boldsymbol{M_{X_{1}}}&\equiv&\boldsymbol{M_{1}}\quad\text{for notational convenience}\\
\boldsymbol{y} &=& \boldsymbol{X_{1}}\beta_{1}+\boldsymbol{X_{2}}\beta_{2}+\boldsymbol{e}\\
\boldsymbol{M_{1}y} &=& \underbrace{\boldsymbol{M_{1}X_{1}\beta_{1}}}_{=0,\text{ since }\boldsymbol{M_{1}X_{1}}=0}+\boldsymbol{M_{1}X_{2}\beta_{2}}+\boldsymbol{\upsilon}\quad\text{(}\boldsymbol{\upsilon}=\boldsymbol{M_{1}e}\text{)}
\end{eqnarray*}
Then: recall that $\boldsymbol{y}=\boldsymbol{P_{x}y}+\boldsymbol{M_{x}y}=\boldsymbol{X_{1}}\beta_{1}+\boldsymbol{X_{2}}\beta_{2}+\boldsymbol{M_{X}y}$\\
multiply previous equation by $\boldsymbol{X'_{2}M_{1}}$, since that
\begin{eqnarray*}
\boldsymbol{X'_{2}M_{1}y}=\overset{\begin{tikzpicture} \node[draw,circle,inner sep=0.4pt]{1};
\end{tikzpicture}}{\boldsymbol{X'_{2}M_{1}X_{1}\hat{\beta_{1}}}}+\overset{\begin{tikzpicture} \node[draw,circle,inner sep=0.4pt]{2};
\end{tikzpicture}}{\boldsymbol{X'_{2}M_{1}X_{2}\hat{\beta_{2}}}}+\overset{\begin{tikzpicture} \node[draw,circle,inner sep=0.4pt]{3};
\end{tikzpicture}}{\boldsymbol{X'_{2}M_{1}M_{X}y}}
\end{eqnarray*}
 of the following properties:
\begin{enumerate}[(1)]
\item $\boldsymbol{M_{1}X_{1}}=0\quad\therefore \boldsymbol{X'_{2}M_{1}X_{1}\beta_{1}}=0$ in \begin{tikzpicture} \node[draw,circle,inner sep=0.4pt]{1};
\end{tikzpicture}
\item $\boldsymbol{M_{1}M_{X}}=\boldsymbol{M_{X}}$ because $M_{X}<M_{1}$\\
$\therefore \boldsymbol{X'_{2}M_{1}M_{X}y}=\boldsymbol{X'_{2}M_{X}y}=\boldsymbol{X'_{2}\hat{e}}=0!$ in \begin{tikzpicture} \node[draw,circle,inner sep=0.4pt]{3};
\end{tikzpicture}
\end{enumerate}
Thus terms (1) and (3) disappear in the above equation and we have
\begin{eqnarray*}
\boldsymbol{X'_{2}M_{1}y} = \boldsymbol{X'_{2}M_{1}X_{2}\hat{\beta_{2}}}
\end{eqnarray*}
Solving for $\boldsymbol{\hat{\beta_{2}}}$, this yields
\begin{eqnarray*}
%I need to use packages to put the equation into textbox
\boldsymbol{\hat{\beta_{2}}}=(\boldsymbol{X'_{2}M_{1}X_{2}})^{-1}\boldsymbol{X'_{2}M_{1}y}
\end{eqnarray*}
$\therefore$ Due to the FWL theorem, we can estimate $\boldsymbol{\beta_{2}}$ directly using orthonogal projections, without having to estimate $\boldsymbol{\beta_{1}}$\\
By extention: $Var(\boldsymbol{\hat{\beta_{2}}})=\sigma^{2}(\boldsymbol{X'_{2}M_{1}X_{2}})^{-1}$

\subsection{Specification issues}
\paragraph{Overspecification} Variables that do not belong to DGP are mistakenly included in DGP.
$\to$ overspecification is not misspecification
\begin{eqnarray*}
\boldsymbol{y}=\boldsymbol{X\beta}+\boldsymbol{Z\gamma}+\boldsymbol{e}\qquad\text{with}\quad \boldsymbol{e}\sim iid(0,\sigma^{2}\boldsymbol{I})\\
\text{but actual DGP is}\quad y=\boldsymbol{y}=\boldsymbol{X\beta_{0}}+\boldsymbol{e}\quad \text{with}\quad \boldsymbol{e}\sim(0,\sigma_{0}^{2}\boldsymbol{I})
\end{eqnarray*}
$\to$ $\beta$ is not biased, since overspecified model is a specific care with
\begin{eqnarray*}
\boldsymbol{\beta}=\boldsymbol{\beta_{0}}\quad\text{and}\quad\boldsymbol{\gamma}=0\quad\text{(plus}\quad\boldsymbol{\sigma^{2}}=\boldsymbol{\sigma_{0}^{2}}\text{)}
\end{eqnarray*}
\paragraph{Underspecification} Important explanatory variables are omitted and we are estimating.
\begin{eqnarray*}
\boldsymbol{y}=\boldsymbol{X\beta}+\boldsymbol{e}\quad \text{with}\quad\boldsymbol{e}\sim(0,\sigma_{0}^{2}\boldsymbol{I})
\end{eqnarray*}
but true DGP is given by
\begin{eqnarray*}
\boldsymbol{y}=\boldsymbol{X\beta_{0}}+\boldsymbol{Z\gamma_{0}}+\boldsymbol{e}\quad \text{with}\quad\boldsymbol{e}\sim(0,\sigma_{0}^{2}\boldsymbol{I})
\end{eqnarray*}
$\to$ \underline{Omitted variable bias}:
\begin{eqnarray*}
\mathbb{E}(\boldsymbol{\hat{\beta}}) &=& \mathbb{E}[(\boldsymbol{X'X})^{-1}\boldsymbol{X'}(\boldsymbol{X\beta_{0}}+\boldsymbol{Z\gamma_{0}}+\boldsymbol{e})]\\
&=& \boldsymbol{\beta_{0}}+(\boldsymbol{X'X})^{-1}\boldsymbol{X'Z\gamma_{0}}+\mathbb{E}[(\boldsymbol{X'X})^{-1}\boldsymbol{X'e}]\\
&=& \boldsymbol{\beta_{0}}+\underbrace{(\boldsymbol{X'X}^{-1}\boldsymbol{X'Z\gamma_{0}})}_{\text{bias}}
\end{eqnarray*}
Whether the bias exists depends on
\begin{enumerate}[(1)]
\item $\boldsymbol{X'Z}=0$ (orthogonal regressor)
\item $\boldsymbol{\gamma}=0$ (model in fact not misspecified)
\end{enumerate}
Comparison between overspecification and underspecification:
\begin{itemize}
\item Overspecification: unbiased, (weakly) consistent, usually insufficient estimates.
\item Underspecification: biased, consistent linear estimator with efficiency that depends on specific sample.
\end{itemize}

\subsection{Heteroskedasticity}
New DGP with $\mathbb{E}(e_{i}^{2}|x_{i})=\sigma^{2}(x_{i})=\sigma^{2}$\\
$y_{i}=\boldsymbol{X_{i}\beta}+e_{i}$\ with $e_{i}\sim id(0,\boldsymbol{\Omega})$\\
$\therefore$ $\boldsymbol{\Omega}\equiv$ diag $(\sigma_{1},\sigma_{2},\cdots,\sigma_{\mu})$\\
OLS is still unbiased when there is heteroskedasticity, because $\mathbb{E}[\beta]$ only depends on first moments and $\mathbb{E}[e]=0$\\
Recall:
\begin{eqnarray*}
\boldsymbol{\hat{\beta}}-\boldsymbol{\beta}=(\boldsymbol{X'X})^{-1}\mathbb{E}[\boldsymbol{X'e}]
\end{eqnarray*}
$\therefore$ The unbiasedness of $\boldsymbol{\hat{\beta}}$ does not depends on higher moments of $\mathbb{E}(e_{i}|x_{i})$
\begin{eqnarray*}
\mathbb{V}(\boldsymbol{\hat{\beta}}) &=& \mathbb{E}[(\boldsymbol{\hat{\beta}}-\boldsymbol{\beta})(\boldsymbol{\hat{\beta}}-\boldsymbol{\beta})']\\
&=& \mathbb{E}[(\boldsymbol{X'X})^{-1}\boldsymbol{X'ee'X}(\boldsymbol{X'X})^{-1}]\\
&=& (\boldsymbol{X'X})^{-1}\boldsymbol{X'}\mathbb{E}[\boldsymbol{ee'}]\boldsymbol{X}(\boldsymbol{X'X})^{-1}\\
&=& (\boldsymbol{X'X})^{-1}(\boldsymbol{X'\Omega X})(\boldsymbol{X'X})^{-1}
\end{eqnarray*}
$\therefore$ GLS estimator: $\boldsymbol{\hat{\beta}}=(\boldsymbol{X'\Omega^{-1}X})^{-1}\boldsymbol{X'\Omega^{-1}y}$\\
$\to$ key challenge when implementing GLS: $\Omega$ is not observable\\
$\therefore$ We must estimate $\boldsymbol{\Omega}$

\subsection{Instrumental variables}
OLS as a general estimator of the form
\begin{eqnarray*}
\boldsymbol{\hat{\beta_{W}}}=(\boldsymbol{W'W})^{-1}\boldsymbol{W'y}\qquad\text{eg.}\quad\boldsymbol{\hat{\beta_{IV}}}\equiv\boldsymbol{\hat{\beta_{2SLS}}}\\
\text{with just-identification(ie. number of parametors is equal to number of instruments)}
\end{eqnarray*}
Two specific cases that can interfere with our key assumption;
\begin{enumerate}[(i)]
\item Measurement error
\item Endogeneity (ie. uncorrelation b/w error and regressors $\quad\mathbb{E}(\boldsymbol{X'e})=0$) due to simutaneous equations. eg. demand \& supply: $q_{i}=-\beta_{1}p_{i}+e_{1i}$, $q_{i}=\beta_{2}p_{i}+e_{2i}$
\end{enumerate}
$\therefore\quad\boldsymbol{\hat{\beta}}$ will be biased and inconsistent!\\
(Angrist \& Krneger, QTE 91)\\
%Picture
\begin{tikzpicture}
    \matrix (m) [matrix of math nodes,row sep=2em,column sep=2em,minimum width=1em]
  {
     & \overset{\text{smoking}}{\boldsymbol{\hat{X}}} & \\
    \overset{\text{month of birth/distance from college}}{\boldsymbol{Z}} & \overset{\text{schooling}}{\boldsymbol{X}} & \overset{\text{wage/health}}{\boldsymbol{y}}\\
   \text{tax rates on tobacco}  & \overset{\text{ability/socioeconomic factors}}{\boldsymbol{e}} & \\};
  \path[-stealth]
    (m-2-1) edge (m-1-2)
            edge (m-2-2)
    (m-2-1) edge [dashed,-] (m-3-2)
    (m-1-2) edge (m-2-3)
    (m-2-2) edge (m-2-3)
    (m-3-2) edge (m-2-3)
    (m-3-1) edge [dashed,-] (m-2-2);
\end{tikzpicture}
$\therefore\quad\frac{dy}{dx}=\beta + \frac{de}{dx}$
\begin{eqnarray*}
\boldsymbol{\hat{\beta_{IV}}}=(\boldsymbol{Z'X})^{-1}\boldsymbol{Z'y}
\end{eqnarray*}
If Z is matrix of k instruments, such that
\begin{eqnarray*}
\boldsymbol{X} &=& [\boldsymbol{X_{1}X_{2}}]\quad\text{where }\boldsymbol{X_{2}}\text{ are endogenous regressors.}\\
\boldsymbol{Z} &=& [\boldsymbol{X_{1}Z_{2}}]\quad\text{where }\boldsymbol{Z_{2}}\text{ are instruments. }\boldsymbol{X_{1}}\text{ are instruments for itself.}
\end{eqnarray*}
Then
\begin{eqnarray*}
\boldsymbol{y}=\boldsymbol{P_{2}X\beta}\quad\text{where}\quad\boldsymbol{P_{2}}=\boldsymbol{Z}(\boldsymbol{Z'Z})^{-1}\boldsymbol{Z'}
\end{eqnarray*}
$\therefore$ 2SLS estimator is IV estimator in just-identified case:
\begin{eqnarray*}
\boldsymbol{\hat{beta_{2SLS}}} &=& (\boldsymbol{X'P_{2}X})^{-1}\boldsymbol{X'P_{2}y}\\
&=& [\underbrace{\boldsymbol{X'Z}(\boldsymbol{Z'Z})^{-1}\boldsymbol{Z'}}_{\boldsymbol{\hat{X}'}}\boldsymbol{X}][\underbrace{\boldsymbol{X'Z}(\boldsymbol{Z'Z})^{-1}\boldsymbol{Z'}}_{\boldsymbol{\hat{X}'}}\boldsymbol{y}]
\end{eqnarray*}
With 1 stage regresson
\begin{eqnarray*}
\boldsymbol{\hat{X}} &=& \boldsymbol{Z}(\boldsymbol{Z'Z})^{-1}\boldsymbol{Z'X}\\
&=& \boldsymbol{Z\hat{\beta}}_{1stage}\quad\text{where}\quad\boldsymbol{\hat{\beta}}_{1stage}=(\boldsymbol{Z'Z})^{-1}\boldsymbol{Z'X}
\end{eqnarray*}
$\to$ estimation is also biased (but generally less so than OLS) and consistent (vs. inconsistent in OLS).
%graph-inconsistency of OLS

\paragraph{Measurement error} ~\\
True regressors $\boldsymbol{X^{*}}$ are not observed and only measured with error $\boldsymbol{\mu}$ (independent of $\boldsymbol{X}$ and $\boldsymbol{y}$)
\begin{eqnarray*}
\boldsymbol{X} &=& \boldsymbol{X^{*}}+\boldsymbol{\mu}\\
\text{True model:}\quad \boldsymbol{y} &=& \boldsymbol{X^{*}\beta}+\boldsymbol{e}\\
&=& \boldsymbol{X}-\boldsymbol{\mu})\boldsymbol{\beta}+\boldsymbol{e}\\
&=& \boldsymbol{X\beta}+\underbrace{\boldsymbol{e}-\boldsymbol{\mu\beta}}_{\nu}\\
&=& \boldsymbol{X\beta}+\boldsymbol{\nu}
\end{eqnarray*}
\begin{eqnarray*}
\to \mathbb{E}[\boldsymbol{\hat{\beta}}] &=& \mathbb{E}[(\boldsymbol{X'X})^{-1}\boldsymbol{X'y}]\\
&=& \mathbb{E}[(\boldsymbol{X'X})^{-1}\boldsymbol{X}(\boldsymbol{X\beta}+\nu)]\\
&=& \boldsymbol{\beta} + \mathbb{E}[(\boldsymbol{X'X})^{-1}\boldsymbol{X'\nu}]
\end{eqnarray*}
$\to$ determine size of bias from $\mathbb{E}[\boldsymbol{X'\nu}]$
\begin{eqnarray*}
\mathbb{E}[\boldsymbol{X'\nu}] &=& \mathbb{E}[(\boldsymbol{X^{*}+\mu})'(\boldsymbol{e-\mu\beta})]\\
&=& \mathbb{E}[\boldsymbol{X^{*'}e}+\boldsymbol{\mu'e}-\boldsymbol{X^{*'}\mu\beta}-\boldsymbol{\mu'\mu\beta}]\\
&=& -\mathbb{E}[\boldsymbol{\mu'\mu}]\boldsymbol{\beta}
\end{eqnarray*}
\begin{eqnarray*}
&& -\mathbb{E}[\boldsymbol{\mu'\mu}]\boldsymbol{\beta}\neq 0\quad\text{if}\quad \boldsymbol{\beta}\neq 0\quad\text{and}\quad\mathbb{E}[\boldsymbol{\mu'\mu}]\neq 0\\
&& \boldsymbol{\hat{\beta}}\xrightarrow{P}\boldsymbol{\beta^{*}}\quad\text{(probability limit of )}\boldsymbol{\hat{\beta}}\\
&& \therefore\quad\boldsymbol{\beta^{*}}=\boldsymbol{\beta}-\mathbb{E}[(\boldsymbol{X'X})^{-1}\boldsymbol{\mu'\mu}]\boldsymbol{\beta}
\end{eqnarray*}

\paragraph{2SLS Computation}
\begin{enumerate}[1)]
    \item Regress $\boldsymbol{X}$ on $\boldsymbol{Z}$  ie $\boldsymbol{\Gamma}=(\boldsymbol{Z'Z})^{-1}\boldsymbol{Z'X}$ and $\boldsymbol{\hat{X}}=\boldsymbol{Z\hat{\Gamma}}$
    \item Regress $\boldsymbol{y}$ on $\boldsymbol{\hat{X}}$    ie $\boldsymbol{\hat{\beta}}=(\boldsymbol{\hat{X'}X'})^{-1}\boldsymbol{\hat{X'}y}$
\end{enumerate}
Scrutinize $\boldsymbol{\hat{X}}$, recalling $\boldsymbol{X}=[\boldsymbol{X_{1}X_{2}}]$, $\boldsymbol{Z}=[\boldsymbol{X_{1}Z_{2}}]$
\begin{eqnarray*}
\boldsymbol{\hat{X}} &=& [\boldsymbol{\hat{X_{1}}},\boldsymbol{\hat{X_{2}}}]\\
&=& [\boldsymbol{P_{2}X_{1}},\boldsymbol{P_{2}X_{2}}]\\
&=& [\boldsymbol{X_{1}},\boldsymbol{P_{2}X_{2}}]\\
&=& [\boldsymbol{X_{1}},\boldsymbol{\hat{X_{2}}}]
\end{eqnarray*}
$\therefore$ Only endogenous variables are replaced with their fitted values
\begin{eqnarray*}
\to\quad\boldsymbol{\hat{X_{2}}}=\boldsymbol{Z_{1}\hat{\Gamma_{12}}}+\boldsymbol{Z_{2}\hat{\Gamma_{22}}}
\end{eqnarray*}